% =====================================================================
%  _ugnn_graph_macros.tex  --  representative graph-signal insets for Fig 1b
%  (reference: oarfish-9)
%
%  Provides \drawugraph{<xc>}{<yc>}{<level>}{<gid>}{<mode>} which renders ONE copy of a
%  FIXED jittered grid graph (same node positions + edges at every level) with:
%    * progressive LOW-PASS node signal   -- deeper levels use a broader bump
%      kernel (smaller decay), i.e. a growing receptive field (GNN convolution).
%    * display MODE (#5)                   -- 0 clean field; 1 noisy encoder feature
%      (clean + LP residual noise nz); 2 clean estimate xhat_0 (clean + small residual
%      nh); 3 noise estimate eps=x_k-xhat_0 (grey 0.5 + LP eps ne). Baked offset tables
%      (nz,ne,nh) come from gen_lpnoise.py -> _ugnn_graph_noise.tex; selection is clean.
%    * nested perturbed SELECTION          -- each node has a baked priority (0..63)
%      from a perturbed rule key=(1-a)*energy+a*uniform (energy-biased but SPREAD,
%      not clustered); active at depth b iff priority < N_b, N_b in {64,32,16,8}
%      (strict x2 down-sampling) -> NESTED subsets (matches D_b = prod_{j<=b} C_j).
%      Priorities are generated by gen_priority.py into _ugnn_graph_priority.tex.
%    * de-highlighted INACTIVE nodes       -- drawn hollow/faint (NaN / zero feat).
%
%  Borrows the multi-radial colour-field idea from diffusion_diagram.tex; the encoder
%  column DISPLAYS a low-pass of the noisy signal x_k (clean + zero-mean node noise),
%  so the residual noise shrinks/smooths with depth (denoising); selection stays clean.
%  The baked LP-noise offsets are generated by gen_lpnoise.py -> _ugnn_graph_noise.tex.
%
%  Node positions depend ONLY on (i,j) -> identical across all calls (the graph is
%  fixed); <level> sets the smoothing + active mask; <gid> is a unique string used
%  only to key the per-node scratch macros (e.g. L1,L2,...,R4).
%
%  REQUIRES tikz + amsmath already loaded.  \input AFTER _ugnn_preamble.tex.
% =====================================================================

% ---- global graph parameters (one source of truth) ------------------
\def\gN{8}                       % grid is gN x gN
\pgfmathtruncatemacro{\gNm}{\gN-1}
\pgfmathtruncatemacro{\gNmm}{\gN-2}
\pgfmathsetmacro{\gC}{(\gN-1)/2}  % centre index (for centring the lattice on 0)
\def\gSp{0.40}                   % grid spacing (cm)
\pgfmathsetmacro{\half}{(\gN-1)/2*\gSp}  % half-extent (for normalising distances)
\def\gJit{0.42}                  % jitter fraction of spacing (organic, graph-like)
\def\gSeed{20260630}             % jitter seed (fixed -> same graph everywhere)
\def\gDbase{2.0}                 % base bump decay (used for the SELECTION score)
\def\gRa{0.072}                  % active node radius (cm)
\def\gRi{0.052}                  % inactive node radius (cm)

% ---- bump base colours (multi-radial signal sources) ----------------
\definecolor{ugA}{RGB}{31,119,180}   % blue
\definecolor{ugB}{RGB}{255,127,14}   % orange
\definecolor{ugC}{RGB}{44,160,44}    % green
\pgfmathsetmacro{\ugAR}{31/255}\pgfmathsetmacro{\ugAG}{119/255}\pgfmathsetmacro{\ugAB}{180/255}
\pgfmathsetmacro{\ugBR}{255/255}\pgfmathsetmacro{\ugBG}{127/255}\pgfmathsetmacro{\ugBB}{14/255}
\pgfmathsetmacro{\ugCR}{44/255}\pgfmathsetmacro{\ugCG}{160/255}\pgfmathsetmacro{\ugCB}{44/255}

% ---- baked nested selection priorities (perturbed energy+uniform; AUTO-GENERATED) ----
% _ugnn_graph_priority.tex  --  AUTO-GENERATED by gen_priority.py (do not hand-edit)
%% nested top-N selection priorities for the fixed 8x8 graph (Fig 1b insets).
%% active at depth b  <=>  pri < N_b, with N_b in {64,32,16,8}.
% rule: key = (1-0.50)*energy_norm + 0.50*uniform(seed=73)  (energy-biased but spread).
\expandafter\def\csname ugP@0@0\endcsname{46}
\expandafter\def\csname ugP@0@1\endcsname{54}
\expandafter\def\csname ugP@0@2\endcsname{56}
\expandafter\def\csname ugP@0@3\endcsname{10}
\expandafter\def\csname ugP@0@4\endcsname{49}
\expandafter\def\csname ugP@0@5\endcsname{59}
\expandafter\def\csname ugP@0@6\endcsname{12}
\expandafter\def\csname ugP@0@7\endcsname{61}
\expandafter\def\csname ugP@1@0\endcsname{18}
\expandafter\def\csname ugP@1@1\endcsname{21}
\expandafter\def\csname ugP@1@2\endcsname{60}
\expandafter\def\csname ugP@1@3\endcsname{57}
\expandafter\def\csname ugP@1@4\endcsname{23}
\expandafter\def\csname ugP@1@5\endcsname{7}
\expandafter\def\csname ugP@1@6\endcsname{47}
\expandafter\def\csname ugP@1@7\endcsname{53}
\expandafter\def\csname ugP@2@0\endcsname{16}
\expandafter\def\csname ugP@2@1\endcsname{37}
\expandafter\def\csname ugP@2@2\endcsname{8}
\expandafter\def\csname ugP@2@3\endcsname{50}
\expandafter\def\csname ugP@2@4\endcsname{0}
\expandafter\def\csname ugP@2@5\endcsname{15}
\expandafter\def\csname ugP@2@6\endcsname{39}
\expandafter\def\csname ugP@2@7\endcsname{58}
\expandafter\def\csname ugP@3@0\endcsname{17}
\expandafter\def\csname ugP@3@1\endcsname{2}
\expandafter\def\csname ugP@3@2\endcsname{19}
\expandafter\def\csname ugP@3@3\endcsname{24}
\expandafter\def\csname ugP@3@4\endcsname{6}
\expandafter\def\csname ugP@3@5\endcsname{41}
\expandafter\def\csname ugP@3@6\endcsname{32}
\expandafter\def\csname ugP@3@7\endcsname{42}
\expandafter\def\csname ugP@4@0\endcsname{38}
\expandafter\def\csname ugP@4@1\endcsname{4}
\expandafter\def\csname ugP@4@2\endcsname{3}
\expandafter\def\csname ugP@4@3\endcsname{5}
\expandafter\def\csname ugP@4@4\endcsname{45}
\expandafter\def\csname ugP@4@5\endcsname{30}
\expandafter\def\csname ugP@4@6\endcsname{20}
\expandafter\def\csname ugP@4@7\endcsname{26}
\expandafter\def\csname ugP@5@0\endcsname{11}
\expandafter\def\csname ugP@5@1\endcsname{25}
\expandafter\def\csname ugP@5@2\endcsname{48}
\expandafter\def\csname ugP@5@3\endcsname{13}
\expandafter\def\csname ugP@5@4\endcsname{14}
\expandafter\def\csname ugP@5@5\endcsname{44}
\expandafter\def\csname ugP@5@6\endcsname{9}
\expandafter\def\csname ugP@5@7\endcsname{43}
\expandafter\def\csname ugP@6@0\endcsname{55}
\expandafter\def\csname ugP@6@1\endcsname{27}
\expandafter\def\csname ugP@6@2\endcsname{36}
\expandafter\def\csname ugP@6@3\endcsname{1}
\expandafter\def\csname ugP@6@4\endcsname{28}
\expandafter\def\csname ugP@6@5\endcsname{35}
\expandafter\def\csname ugP@6@6\endcsname{51}
\expandafter\def\csname ugP@6@7\endcsname{33}
\expandafter\def\csname ugP@7@0\endcsname{63}
\expandafter\def\csname ugP@7@1\endcsname{22}
\expandafter\def\csname ugP@7@2\endcsname{29}
\expandafter\def\csname ugP@7@3\endcsname{34}
\expandafter\def\csname ugP@7@4\endcsname{31}
\expandafter\def\csname ugP@7@5\endcsname{62}
\expandafter\def\csname ugP@7@6\endcsname{40}
\expandafter\def\csname ugP@7@7\endcsname{52}

% ---- baked low-pass noise offsets (forward x_k corruption, denoised by depth; AUTO-GEN) ----
% _ugnn_graph_noise.tex  --  AUTO-GENERATED by gen_lpnoise.py (do not hand-edit)
%% nz = encoder residual noise (LP of x_k corruption);  base = clean_smoothed (mode 1)
%% ne = decoder noise estimate eps = x_k - xhat_0 (LP);  base = grey 0.5      (mode 3)
%% nh = xhat_0 small residual error (full-res);          base = clean         (mode 2)
%% all zero-mean, renormalised to per-depth RMS floors so variability survives.
% AMP_BIG=0.60 AMP_SMALL=0.18 SIGMA={1: 0.0, 2: 0.45, 3: 0.7, 4: 0.95} ENC_RMS={1: 0.19, 2: 0.16, 3: 0.14, 4: 0.12} EPS_RMS={1: 0.185, 2: 0.16, 3: 0.14, 4: 0.12} HAT_RMS=0.085 SEED=2027
\expandafter\def\csname nzR@1@0@0\endcsname{-0.3140}
\expandafter\def\csname nzG@1@0@0\endcsname{-0.0728}
\expandafter\def\csname nzB@1@0@0\endcsname{-0.2664}
\expandafter\def\csname neR@1@0@0\endcsname{-0.2688}
\expandafter\def\csname neG@1@0@0\endcsname{-0.0045}
\expandafter\def\csname neB@1@0@0\endcsname{-0.2764}
\expandafter\def\csname nzR@1@0@1\endcsname{-0.0013}
\expandafter\def\csname nzG@1@0@1\endcsname{-0.0327}
\expandafter\def\csname nzB@1@0@1\endcsname{0.1311}
\expandafter\def\csname neR@1@0@1\endcsname{-0.0549}
\expandafter\def\csname neG@1@0@1\endcsname{0.0296}
\expandafter\def\csname neB@1@0@1\endcsname{0.1390}
\expandafter\def\csname nzR@1@0@2\endcsname{-0.1291}
\expandafter\def\csname nzG@1@0@2\endcsname{0.0774}
\expandafter\def\csname nzB@1@0@2\endcsname{0.0389}
\expandafter\def\csname neR@1@0@2\endcsname{-0.1471}
\expandafter\def\csname neG@1@0@2\endcsname{0.1374}
\expandafter\def\csname neB@1@0@2\endcsname{-0.0140}
\expandafter\def\csname nzR@1@0@3\endcsname{-0.2618}
\expandafter\def\csname nzG@1@0@3\endcsname{-0.3134}
\expandafter\def\csname nzB@1@0@3\endcsname{-0.0751}
\expandafter\def\csname neR@1@0@3\endcsname{-0.1635}
\expandafter\def\csname neG@1@0@3\endcsname{-0.3095}
\expandafter\def\csname neB@1@0@3\endcsname{-0.0061}
\expandafter\def\csname nzR@1@0@4\endcsname{0.2533}
\expandafter\def\csname nzG@1@0@4\endcsname{-0.3006}
\expandafter\def\csname nzB@1@0@4\endcsname{0.2945}
\expandafter\def\csname neR@1@0@4\endcsname{0.1685}
\expandafter\def\csname neG@1@0@4\endcsname{-0.2376}
\expandafter\def\csname neB@1@0@4\endcsname{0.3076}
\expandafter\def\csname nzR@1@0@5\endcsname{0.1565}
\expandafter\def\csname nzG@1@0@5\endcsname{0.1427}
\expandafter\def\csname nzB@1@0@5\endcsname{-0.2019}
\expandafter\def\csname neR@1@0@5\endcsname{0.1444}
\expandafter\def\csname neG@1@0@5\endcsname{0.0953}
\expandafter\def\csname neB@1@0@5\endcsname{-0.1713}
\expandafter\def\csname nzR@1@0@6\endcsname{0.0556}
\expandafter\def\csname nzG@1@0@6\endcsname{0.1381}
\expandafter\def\csname nzB@1@0@6\endcsname{0.2611}
\expandafter\def\csname neR@1@0@6\endcsname{-0.0370}
\expandafter\def\csname neG@1@0@6\endcsname{0.2022}
\expandafter\def\csname neB@1@0@6\endcsname{0.2081}
\expandafter\def\csname nzR@1@0@7\endcsname{-0.2928}
\expandafter\def\csname nzG@1@0@7\endcsname{-0.2541}
\expandafter\def\csname nzB@1@0@7\endcsname{0.2998}
\expandafter\def\csname neR@1@0@7\endcsname{-0.3034}
\expandafter\def\csname neG@1@0@7\endcsname{-0.1596}
\expandafter\def\csname neB@1@0@7\endcsname{0.2262}
\expandafter\def\csname nzR@1@1@0\endcsname{0.1890}
\expandafter\def\csname nzG@1@1@0\endcsname{-0.0901}
\expandafter\def\csname nzB@1@1@0\endcsname{-0.2330}
\expandafter\def\csname neR@1@1@0\endcsname{0.1337}
\expandafter\def\csname neG@1@1@0\endcsname{-0.0564}
\expandafter\def\csname neB@1@1@0\endcsname{-0.2745}
\expandafter\def\csname nzR@1@1@1\endcsname{0.0938}
\expandafter\def\csname nzG@1@1@1\endcsname{0.2694}
\expandafter\def\csname nzB@1@1@1\endcsname{-0.1858}
\expandafter\def\csname neR@1@1@1\endcsname{0.0380}
\expandafter\def\csname neG@1@1@1\endcsname{0.2243}
\expandafter\def\csname neB@1@1@1\endcsname{-0.2260}
\expandafter\def\csname nzR@1@1@2\endcsname{0.1449}
\expandafter\def\csname nzG@1@1@2\endcsname{-0.2014}
\expandafter\def\csname nzB@1@1@2\endcsname{0.1428}
\expandafter\def\csname neR@1@1@2\endcsname{0.1705}
\expandafter\def\csname neG@1@1@2\endcsname{-0.2038}
\expandafter\def\csname neB@1@1@2\endcsname{0.2010}
\expandafter\def\csname nzR@1@1@3\endcsname{-0.2299}
\expandafter\def\csname nzG@1@1@3\endcsname{0.3136}
\expandafter\def\csname nzB@1@1@3\endcsname{-0.2879}
\expandafter\def\csname neR@1@1@3\endcsname{-0.1935}
\expandafter\def\csname neG@1@1@3\endcsname{0.3719}
\expandafter\def\csname neB@1@1@3\endcsname{-0.2900}
\expandafter\def\csname nzR@1@1@4\endcsname{-0.0909}
\expandafter\def\csname nzG@1@1@4\endcsname{0.0901}
\expandafter\def\csname nzB@1@1@4\endcsname{0.0774}
\expandafter\def\csname neR@1@1@4\endcsname{-0.0823}
\expandafter\def\csname neG@1@1@4\endcsname{0.0052}
\expandafter\def\csname neB@1@1@4\endcsname{-0.0119}
\expandafter\def\csname nzR@1@1@5\endcsname{-0.2102}
\expandafter\def\csname nzG@1@1@5\endcsname{-0.2557}
\expandafter\def\csname nzB@1@1@5\endcsname{-0.1254}
\expandafter\def\csname neR@1@1@5\endcsname{-0.2667}
\expandafter\def\csname neG@1@1@5\endcsname{-0.1839}
\expandafter\def\csname neB@1@1@5\endcsname{-0.0924}
\expandafter\def\csname nzR@1@1@6\endcsname{-0.1786}
\expandafter\def\csname nzG@1@1@6\endcsname{0.1513}
\expandafter\def\csname nzB@1@1@6\endcsname{-0.1016}
\expandafter\def\csname neR@1@1@6\endcsname{-0.1877}
\expandafter\def\csname neG@1@1@6\endcsname{0.1139}
\expandafter\def\csname neB@1@1@6\endcsname{-0.1089}
\expandafter\def\csname nzR@1@1@7\endcsname{0.1171}
\expandafter\def\csname nzG@1@1@7\endcsname{0.3019}
\expandafter\def\csname nzB@1@1@7\endcsname{0.1226}
\expandafter\def\csname neR@1@1@7\endcsname{0.1408}
\expandafter\def\csname neG@1@1@7\endcsname{0.2728}
\expandafter\def\csname neB@1@1@7\endcsname{0.1756}
\expandafter\def\csname nzR@1@2@0\endcsname{-0.0934}
\expandafter\def\csname nzG@1@2@0\endcsname{0.1239}
\expandafter\def\csname nzB@1@2@0\endcsname{0.2489}
\expandafter\def\csname neR@1@2@0\endcsname{-0.0637}
\expandafter\def\csname neG@1@2@0\endcsname{0.1054}
\expandafter\def\csname neB@1@2@0\endcsname{0.1626}
\expandafter\def\csname nzR@1@2@1\endcsname{-0.0184}
\expandafter\def\csname nzG@1@2@1\endcsname{0.0484}
\expandafter\def\csname nzB@1@2@1\endcsname{-0.1371}
\expandafter\def\csname neR@1@2@1\endcsname{-0.0870}
\expandafter\def\csname neG@1@2@1\endcsname{0.1036}
\expandafter\def\csname neB@1@2@1\endcsname{-0.0398}
\expandafter\def\csname nzR@1@2@2\endcsname{0.1083}
\expandafter\def\csname nzG@1@2@2\endcsname{-0.1882}
\expandafter\def\csname nzB@1@2@2\endcsname{0.1443}
\expandafter\def\csname neR@1@2@2\endcsname{0.0269}
\expandafter\def\csname neG@1@2@2\endcsname{-0.1704}
\expandafter\def\csname neB@1@2@2\endcsname{0.1908}
\expandafter\def\csname nzR@1@2@3\endcsname{0.2849}
\expandafter\def\csname nzG@1@2@3\endcsname{0.1501}
\expandafter\def\csname nzB@1@2@3\endcsname{-0.2144}
\expandafter\def\csname neR@1@2@3\endcsname{0.3342}
\expandafter\def\csname neG@1@2@3\endcsname{0.2085}
\expandafter\def\csname neB@1@2@3\endcsname{-0.1774}
\expandafter\def\csname nzR@1@2@4\endcsname{-0.2195}
\expandafter\def\csname nzG@1@2@4\endcsname{0.2475}
\expandafter\def\csname nzB@1@2@4\endcsname{0.0920}
\expandafter\def\csname neR@1@2@4\endcsname{-0.2414}
\expandafter\def\csname neG@1@2@4\endcsname{0.1900}
\expandafter\def\csname neB@1@2@4\endcsname{0.0225}
\expandafter\def\csname nzR@1@2@5\endcsname{-0.2527}
\expandafter\def\csname nzG@1@2@5\endcsname{0.2266}
\expandafter\def\csname nzB@1@2@5\endcsname{0.0437}
\expandafter\def\csname neR@1@2@5\endcsname{-0.1781}
\expandafter\def\csname neG@1@2@5\endcsname{0.1508}
\expandafter\def\csname neB@1@2@5\endcsname{0.1026}
\expandafter\def\csname nzR@1@2@6\endcsname{-0.0554}
\expandafter\def\csname nzG@1@2@6\endcsname{0.2033}
\expandafter\def\csname nzB@1@2@6\endcsname{0.0100}
\expandafter\def\csname neR@1@2@6\endcsname{-0.0044}
\expandafter\def\csname neG@1@2@6\endcsname{0.2622}
\expandafter\def\csname neB@1@2@6\endcsname{-0.0678}
\expandafter\def\csname nzR@1@2@7\endcsname{0.1590}
\expandafter\def\csname nzG@1@2@7\endcsname{-0.1170}
\expandafter\def\csname nzB@1@2@7\endcsname{-0.3035}
\expandafter\def\csname neR@1@2@7\endcsname{0.1393}
\expandafter\def\csname neG@1@2@7\endcsname{-0.0996}
\expandafter\def\csname neB@1@2@7\endcsname{-0.3048}
\expandafter\def\csname nzR@1@3@0\endcsname{-0.1777}
\expandafter\def\csname nzG@1@3@0\endcsname{0.1116}
\expandafter\def\csname nzB@1@3@0\endcsname{-0.0033}
\expandafter\def\csname neR@1@3@0\endcsname{-0.2285}
\expandafter\def\csname neG@1@3@0\endcsname{0.0875}
\expandafter\def\csname neB@1@3@0\endcsname{0.0678}
\expandafter\def\csname nzR@1@3@1\endcsname{0.1201}
\expandafter\def\csname nzG@1@3@1\endcsname{0.0571}
\expandafter\def\csname nzB@1@3@1\endcsname{0.1110}
\expandafter\def\csname neR@1@3@1\endcsname{0.0933}
\expandafter\def\csname neG@1@3@1\endcsname{-0.0309}
\expandafter\def\csname neB@1@3@1\endcsname{0.0640}
\expandafter\def\csname nzR@1@3@2\endcsname{-0.0828}
\expandafter\def\csname nzG@1@3@2\endcsname{-0.1538}
\expandafter\def\csname nzB@1@3@2\endcsname{0.1155}
\expandafter\def\csname neR@1@3@2\endcsname{-0.1551}
\expandafter\def\csname neG@1@3@2\endcsname{-0.1607}
\expandafter\def\csname neB@1@3@2\endcsname{0.1363}
\expandafter\def\csname nzR@1@3@3\endcsname{0.0947}
\expandafter\def\csname nzG@1@3@3\endcsname{-0.3111}
\expandafter\def\csname nzB@1@3@3\endcsname{0.0612}
\expandafter\def\csname neR@1@3@3\endcsname{0.1641}
\expandafter\def\csname neG@1@3@3\endcsname{-0.3520}
\expandafter\def\csname neB@1@3@3\endcsname{-0.0284}
\expandafter\def\csname nzR@1@3@4\endcsname{-0.1352}
\expandafter\def\csname nzG@1@3@4\endcsname{0.2930}
\expandafter\def\csname nzB@1@3@4\endcsname{-0.2243}
\expandafter\def\csname neR@1@3@4\endcsname{-0.0834}
\expandafter\def\csname neG@1@3@4\endcsname{0.3288}
\expandafter\def\csname neB@1@3@4\endcsname{-0.1926}
\expandafter\def\csname nzR@1@3@5\endcsname{0.0126}
\expandafter\def\csname nzG@1@3@5\endcsname{-0.0898}
\expandafter\def\csname nzB@1@3@5\endcsname{0.0309}
\expandafter\def\csname neR@1@3@5\endcsname{0.0921}
\expandafter\def\csname neG@1@3@5\endcsname{-0.0520}
\expandafter\def\csname neB@1@3@5\endcsname{0.0124}
\expandafter\def\csname nzR@1@3@6\endcsname{0.1593}
\expandafter\def\csname nzG@1@3@6\endcsname{0.1897}
\expandafter\def\csname nzB@1@3@6\endcsname{-0.1561}
\expandafter\def\csname neR@1@3@6\endcsname{0.1709}
\expandafter\def\csname neG@1@3@6\endcsname{0.1495}
\expandafter\def\csname neB@1@3@6\endcsname{-0.1427}
\expandafter\def\csname nzR@1@3@7\endcsname{0.0522}
\expandafter\def\csname nzG@1@3@7\endcsname{0.1370}
\expandafter\def\csname nzB@1@3@7\endcsname{0.1081}
\expandafter\def\csname neR@1@3@7\endcsname{0.0249}
\expandafter\def\csname neG@1@3@7\endcsname{0.2046}
\expandafter\def\csname neB@1@3@7\endcsname{0.0938}
\expandafter\def\csname nzR@1@4@0\endcsname{0.2904}
\expandafter\def\csname nzG@1@4@0\endcsname{0.0547}
\expandafter\def\csname nzB@1@4@0\endcsname{0.1143}
\expandafter\def\csname neR@1@4@0\endcsname{0.2272}
\expandafter\def\csname neG@1@4@0\endcsname{0.0571}
\expandafter\def\csname neB@1@4@0\endcsname{0.1489}
\expandafter\def\csname nzR@1@4@1\endcsname{-0.3087}
\expandafter\def\csname nzG@1@4@1\endcsname{-0.2245}
\expandafter\def\csname nzB@1@4@1\endcsname{-0.1819}
\expandafter\def\csname neR@1@4@1\endcsname{-0.3091}
\expandafter\def\csname neG@1@4@1\endcsname{-0.2914}
\expandafter\def\csname neB@1@4@1\endcsname{-0.1098}
\expandafter\def\csname nzR@1@4@2\endcsname{-0.0804}
\expandafter\def\csname nzG@1@4@2\endcsname{-0.1802}
\expandafter\def\csname nzB@1@4@2\endcsname{0.0549}
\expandafter\def\csname neR@1@4@2\endcsname{-0.1271}
\expandafter\def\csname neG@1@4@2\endcsname{-0.2490}
\expandafter\def\csname neB@1@4@2\endcsname{0.1237}
\expandafter\def\csname nzR@1@4@3\endcsname{-0.2205}
\expandafter\def\csname nzG@1@4@3\endcsname{0.0242}
\expandafter\def\csname nzB@1@4@3\endcsname{-0.1754}
\expandafter\def\csname neR@1@4@3\endcsname{-0.2597}
\expandafter\def\csname neG@1@4@3\endcsname{-0.0347}
\expandafter\def\csname neB@1@4@3\endcsname{-0.1228}
\expandafter\def\csname nzR@1@4@4\endcsname{0.2242}
\expandafter\def\csname nzG@1@4@4\endcsname{-0.3046}
\expandafter\def\csname nzB@1@4@4\endcsname{-0.2986}
\expandafter\def\csname neR@1@4@4\endcsname{0.1916}
\expandafter\def\csname neG@1@4@4\endcsname{-0.2967}
\expandafter\def\csname neB@1@4@4\endcsname{-0.3354}
\expandafter\def\csname nzR@1@4@5\endcsname{0.0113}
\expandafter\def\csname nzG@1@4@5\endcsname{-0.2473}
\expandafter\def\csname nzB@1@4@5\endcsname{0.0106}
\expandafter\def\csname neR@1@4@5\endcsname{0.0018}
\expandafter\def\csname neG@1@4@5\endcsname{-0.1948}
\expandafter\def\csname neB@1@4@5\endcsname{-0.0336}
\expandafter\def\csname nzR@1@4@6\endcsname{-0.2623}
\expandafter\def\csname nzG@1@4@6\endcsname{0.1949}
\expandafter\def\csname nzB@1@4@6\endcsname{0.1855}
\expandafter\def\csname neR@1@4@6\endcsname{-0.3089}
\expandafter\def\csname neG@1@4@6\endcsname{0.2194}
\expandafter\def\csname neB@1@4@6\endcsname{0.2118}
\expandafter\def\csname nzR@1@4@7\endcsname{-0.0645}
\expandafter\def\csname nzG@1@4@7\endcsname{-0.1880}
\expandafter\def\csname nzB@1@4@7\endcsname{-0.1184}
\expandafter\def\csname neR@1@4@7\endcsname{-0.0327}
\expandafter\def\csname neG@1@4@7\endcsname{-0.1283}
\expandafter\def\csname neB@1@4@7\endcsname{-0.1270}
\expandafter\def\csname nzR@1@5@0\endcsname{0.2077}
\expandafter\def\csname nzG@1@5@0\endcsname{-0.1958}
\expandafter\def\csname nzB@1@5@0\endcsname{0.2869}
\expandafter\def\csname neR@1@5@0\endcsname{0.1684}
\expandafter\def\csname neG@1@5@0\endcsname{-0.2591}
\expandafter\def\csname neB@1@5@0\endcsname{0.2348}
\expandafter\def\csname nzR@1@5@1\endcsname{0.0642}
\expandafter\def\csname nzG@1@5@1\endcsname{0.2809}
\expandafter\def\csname nzB@1@5@1\endcsname{-0.2379}
\expandafter\def\csname neR@1@5@1\endcsname{0.0164}
\expandafter\def\csname neG@1@5@1\endcsname{0.1956}
\expandafter\def\csname neB@1@5@1\endcsname{-0.2156}
\expandafter\def\csname nzR@1@5@2\endcsname{-0.2743}
\expandafter\def\csname nzG@1@5@2\endcsname{-0.2747}
\expandafter\def\csname nzB@1@5@2\endcsname{-0.2181}
\expandafter\def\csname neR@1@5@2\endcsname{-0.3115}
\expandafter\def\csname neG@1@5@2\endcsname{-0.2421}
\expandafter\def\csname neB@1@5@2\endcsname{-0.1838}
\expandafter\def\csname nzR@1@5@3\endcsname{-0.2364}
\expandafter\def\csname nzG@1@5@3\endcsname{0.2942}
\expandafter\def\csname nzB@1@5@3\endcsname{0.2582}
\expandafter\def\csname neR@1@5@3\endcsname{-0.1964}
\expandafter\def\csname neG@1@5@3\endcsname{0.2443}
\expandafter\def\csname neB@1@5@3\endcsname{0.2051}
\expandafter\def\csname nzR@1@5@4\endcsname{-0.0632}
\expandafter\def\csname nzG@1@5@4\endcsname{0.0727}
\expandafter\def\csname nzB@1@5@4\endcsname{-0.2339}
\expandafter\def\csname neR@1@5@4\endcsname{-0.1116}
\expandafter\def\csname neG@1@5@4\endcsname{0.0049}
\expandafter\def\csname neB@1@5@4\endcsname{-0.2730}
\expandafter\def\csname nzR@1@5@5\endcsname{0.2804}
\expandafter\def\csname nzG@1@5@5\endcsname{-0.3011}
\expandafter\def\csname nzB@1@5@5\endcsname{0.1447}
\expandafter\def\csname neR@1@5@5\endcsname{0.2008}
\expandafter\def\csname neG@1@5@5\endcsname{-0.3350}
\expandafter\def\csname neB@1@5@5\endcsname{0.1565}
\expandafter\def\csname nzR@1@5@6\endcsname{-0.0174}
\expandafter\def\csname nzG@1@5@6\endcsname{-0.2793}
\expandafter\def\csname nzB@1@5@6\endcsname{0.3064}
\expandafter\def\csname neR@1@5@6\endcsname{-0.0798}
\expandafter\def\csname neG@1@5@6\endcsname{-0.3299}
\expandafter\def\csname neB@1@5@6\endcsname{0.2932}
\expandafter\def\csname nzR@1@5@7\endcsname{-0.1192}
\expandafter\def\csname nzG@1@5@7\endcsname{-0.2384}
\expandafter\def\csname nzB@1@5@7\endcsname{0.0683}
\expandafter\def\csname neR@1@5@7\endcsname{-0.1310}
\expandafter\def\csname neG@1@5@7\endcsname{-0.2788}
\expandafter\def\csname neB@1@5@7\endcsname{0.1259}
\expandafter\def\csname nzR@1@6@0\endcsname{-0.3002}
\expandafter\def\csname nzG@1@6@0\endcsname{-0.1046}
\expandafter\def\csname nzB@1@6@0\endcsname{0.2941}
\expandafter\def\csname neR@1@6@0\endcsname{-0.2374}
\expandafter\def\csname neG@1@6@0\endcsname{-0.1579}
\expandafter\def\csname neB@1@6@0\endcsname{0.2843}
\expandafter\def\csname nzR@1@6@1\endcsname{0.2405}
\expandafter\def\csname nzG@1@6@1\endcsname{-0.1658}
\expandafter\def\csname nzB@1@6@1\endcsname{0.0369}
\expandafter\def\csname neR@1@6@1\endcsname{0.1501}
\expandafter\def\csname neG@1@6@1\endcsname{-0.2008}
\expandafter\def\csname neB@1@6@1\endcsname{0.0024}
\expandafter\def\csname nzR@1@6@2\endcsname{0.1853}
\expandafter\def\csname nzG@1@6@2\endcsname{0.3120}
\expandafter\def\csname nzB@1@6@2\endcsname{-0.1886}
\expandafter\def\csname neR@1@6@2\endcsname{0.1503}
\expandafter\def\csname neG@1@6@2\endcsname{0.3085}
\expandafter\def\csname neB@1@6@2\endcsname{-0.1073}
\expandafter\def\csname nzR@1@6@3\endcsname{0.0105}
\expandafter\def\csname nzG@1@6@3\endcsname{-0.0565}
\expandafter\def\csname nzB@1@6@3\endcsname{0.0704}
\expandafter\def\csname neR@1@6@3\endcsname{0.0607}
\expandafter\def\csname neG@1@6@3\endcsname{0.0177}
\expandafter\def\csname neB@1@6@3\endcsname{0.0665}
\expandafter\def\csname nzR@1@6@4\endcsname{-0.2550}
\expandafter\def\csname nzG@1@6@4\endcsname{-0.1660}
\expandafter\def\csname nzB@1@6@4\endcsname{-0.2631}
\expandafter\def\csname neR@1@6@4\endcsname{-0.2568}
\expandafter\def\csname neG@1@6@4\endcsname{-0.2264}
\expandafter\def\csname neB@1@6@4\endcsname{-0.1688}
\expandafter\def\csname nzR@1@6@5\endcsname{0.1063}
\expandafter\def\csname nzG@1@6@5\endcsname{-0.0257}
\expandafter\def\csname nzB@1@6@5\endcsname{0.1842}
\expandafter\def\csname neR@1@6@5\endcsname{0.1540}
\expandafter\def\csname neG@1@6@5\endcsname{-0.0458}
\expandafter\def\csname neB@1@6@5\endcsname{0.1049}
\expandafter\def\csname nzR@1@6@6\endcsname{0.1686}
\expandafter\def\csname nzG@1@6@6\endcsname{0.1890}
\expandafter\def\csname nzB@1@6@6\endcsname{-0.1022}
\expandafter\def\csname neR@1@6@6\endcsname{0.1667}
\expandafter\def\csname neG@1@6@6\endcsname{0.1541}
\expandafter\def\csname neB@1@6@6\endcsname{-0.0167}
\expandafter\def\csname nzR@1@6@7\endcsname{-0.2958}
\expandafter\def\csname nzG@1@6@7\endcsname{0.2596}
\expandafter\def\csname nzB@1@6@7\endcsname{-0.0527}
\expandafter\def\csname neR@1@6@7\endcsname{-0.3550}
\expandafter\def\csname neG@1@6@7\endcsname{0.3107}
\expandafter\def\csname neB@1@6@7\endcsname{-0.0955}
\expandafter\def\csname nzR@1@7@0\endcsname{0.2733}
\expandafter\def\csname nzG@1@7@0\endcsname{-0.0038}
\expandafter\def\csname nzB@1@7@0\endcsname{-0.0652}
\expandafter\def\csname neR@1@7@0\endcsname{0.1792}
\expandafter\def\csname neG@1@7@0\endcsname{-0.0520}
\expandafter\def\csname neB@1@7@0\endcsname{-0.1026}
\expandafter\def\csname nzR@1@7@1\endcsname{0.2150}
\expandafter\def\csname nzG@1@7@1\endcsname{0.1695}
\expandafter\def\csname nzB@1@7@1\endcsname{-0.1616}
\expandafter\def\csname neR@1@7@1\endcsname{0.1893}
\expandafter\def\csname neG@1@7@1\endcsname{0.1642}
\expandafter\def\csname neB@1@7@1\endcsname{-0.1870}
\expandafter\def\csname nzR@1@7@2\endcsname{-0.2456}
\expandafter\def\csname nzG@1@7@2\endcsname{-0.0301}
\expandafter\def\csname nzB@1@7@2\endcsname{-0.0379}
\expandafter\def\csname neR@1@7@2\endcsname{-0.1679}
\expandafter\def\csname neG@1@7@2\endcsname{0.0233}
\expandafter\def\csname neB@1@7@2\endcsname{-0.0480}
\expandafter\def\csname nzR@1@7@3\endcsname{-0.2002}
\expandafter\def\csname nzG@1@7@3\endcsname{-0.0319}
\expandafter\def\csname nzB@1@7@3\endcsname{0.2036}
\expandafter\def\csname neR@1@7@3\endcsname{-0.1627}
\expandafter\def\csname neG@1@7@3\endcsname{0.0443}
\expandafter\def\csname neB@1@7@3\endcsname{0.2127}
\expandafter\def\csname nzR@1@7@4\endcsname{0.2989}
\expandafter\def\csname nzG@1@7@4\endcsname{0.0829}
\expandafter\def\csname nzB@1@7@4\endcsname{-0.0188}
\expandafter\def\csname neR@1@7@4\endcsname{0.2700}
\expandafter\def\csname neG@1@7@4\endcsname{0.1637}
\expandafter\def\csname neB@1@7@4\endcsname{-0.0764}
\expandafter\def\csname nzR@1@7@5\endcsname{-0.0187}
\expandafter\def\csname nzG@1@7@5\endcsname{-0.2352}
\expandafter\def\csname nzB@1@7@5\endcsname{0.2286}
\expandafter\def\csname neR@1@7@5\endcsname{0.0281}
\expandafter\def\csname neG@1@7@5\endcsname{-0.1648}
\expandafter\def\csname neB@1@7@5\endcsname{0.2839}
\expandafter\def\csname nzR@1@7@6\endcsname{-0.1622}
\expandafter\def\csname nzG@1@7@6\endcsname{0.1262}
\expandafter\def\csname nzB@1@7@6\endcsname{-0.1336}
\expandafter\def\csname neR@1@7@6\endcsname{-0.2132}
\expandafter\def\csname neG@1@7@6\endcsname{0.0507}
\expandafter\def\csname neB@1@7@6\endcsname{-0.1227}
\expandafter\def\csname nzR@1@7@7\endcsname{0.2706}
\expandafter\def\csname nzG@1@7@7\endcsname{0.2555}
\expandafter\def\csname nzB@1@7@7\endcsname{0.2352}
\expandafter\def\csname neR@1@7@7\endcsname{0.2119}
\expandafter\def\csname neG@1@7@7\endcsname{0.3082}
\expandafter\def\csname neB@1@7@7\endcsname{0.1802}
\expandafter\def\csname nzR@2@0@0\endcsname{-0.2839}
\expandafter\def\csname nzG@2@0@0\endcsname{-0.0776}
\expandafter\def\csname nzB@2@0@0\endcsname{-0.2638}
\expandafter\def\csname neR@2@0@0\endcsname{-0.2538}
\expandafter\def\csname neG@2@0@0\endcsname{-0.0049}
\expandafter\def\csname neB@2@0@0\endcsname{-0.2806}
\expandafter\def\csname nzR@2@0@1\endcsname{-0.0253}
\expandafter\def\csname nzG@2@0@1\endcsname{-0.0103}
\expandafter\def\csname nzB@2@0@1\endcsname{0.0845}
\expandafter\def\csname neR@2@0@1\endcsname{-0.0762}
\expandafter\def\csname neG@2@0@1\endcsname{0.0522}
\expandafter\def\csname neB@2@0@1\endcsname{0.0852}
\expandafter\def\csname nzR@2@0@2\endcsname{-0.1240}
\expandafter\def\csname nzG@2@0@2\endcsname{0.0312}
\expandafter\def\csname nzB@2@0@2\endcsname{0.0463}
\expandafter\def\csname neR@2@0@2\endcsname{-0.1369}
\expandafter\def\csname neG@2@0@2\endcsname{0.0906}
\expandafter\def\csname neB@2@0@2\endcsname{0.0095}
\expandafter\def\csname nzR@2@0@3\endcsname{-0.2393}
\expandafter\def\csname nzG@2@0@3\endcsname{-0.2714}
\expandafter\def\csname nzB@2@0@3\endcsname{-0.0617}
\expandafter\def\csname neR@2@0@3\endcsname{-0.1595}
\expandafter\def\csname neG@2@0@3\endcsname{-0.2588}
\expandafter\def\csname neB@2@0@3\endcsname{-0.0040}
\expandafter\def\csname nzR@2@0@4\endcsname{0.2068}
\expandafter\def\csname nzG@2@0@4\endcsname{-0.2718}
\expandafter\def\csname nzB@2@0@4\endcsname{0.2432}
\expandafter\def\csname neR@2@0@4\endcsname{0.1408}
\expandafter\def\csname neG@2@0@4\endcsname{-0.2283}
\expandafter\def\csname neB@2@0@4\endcsname{0.2597}
\expandafter\def\csname nzR@2@0@5\endcsname{0.1442}
\expandafter\def\csname nzG@2@0@5\endcsname{0.0965}
\expandafter\def\csname nzB@2@0@5\endcsname{-0.1466}
\expandafter\def\csname neR@2@0@5\endcsname{0.1178}
\expandafter\def\csname neG@2@0@5\endcsname{0.0697}
\expandafter\def\csname neB@2@0@5\endcsname{-0.1226}
\expandafter\def\csname nzR@2@0@6\endcsname{0.0251}
\expandafter\def\csname nzG@2@0@6\endcsname{0.1255}
\expandafter\def\csname nzB@2@0@6\endcsname{0.2309}
\expandafter\def\csname neR@2@0@6\endcsname{-0.0605}
\expandafter\def\csname neG@2@0@6\endcsname{0.1861}
\expandafter\def\csname neB@2@0@6\endcsname{0.1835}
\expandafter\def\csname nzR@2@0@7\endcsname{-0.2667}
\expandafter\def\csname nzG@2@0@7\endcsname{-0.2059}
\expandafter\def\csname nzB@2@0@7\endcsname{0.3164}
\expandafter\def\csname neR@2@0@7\endcsname{-0.2868}
\expandafter\def\csname neG@2@0@7\endcsname{-0.1150}
\expandafter\def\csname neB@2@0@7\endcsname{0.2499}
\expandafter\def\csname nzR@2@1@0\endcsname{0.1437}
\expandafter\def\csname nzG@2@1@0\endcsname{-0.0556}
\expandafter\def\csname nzB@2@1@0\endcsname{-0.2215}
\expandafter\def\csname neR@2@1@0\endcsname{0.0968}
\expandafter\def\csname neG@2@1@0\endcsname{-0.0251}
\expandafter\def\csname neB@2@1@0\endcsname{-0.2720}
\expandafter\def\csname nzR@2@1@1\endcsname{0.0963}
\expandafter\def\csname nzG@2@1@1\endcsname{0.2017}
\expandafter\def\csname nzB@2@1@1\endcsname{-0.1583}
\expandafter\def\csname neR@2@1@1\endcsname{0.0404}
\expandafter\def\csname neG@2@1@1\endcsname{0.1785}
\expandafter\def\csname neB@2@1@1\endcsname{-0.1863}
\expandafter\def\csname nzR@2@1@2\endcsname{0.1081}
\expandafter\def\csname nzG@2@1@2\endcsname{-0.1335}
\expandafter\def\csname nzB@2@1@2\endcsname{0.0953}
\expandafter\def\csname neR@2@1@2\endcsname{0.1229}
\expandafter\def\csname neG@2@1@2\endcsname{-0.1300}
\expandafter\def\csname neB@2@1@2\endcsname{0.1432}
\expandafter\def\csname nzR@2@1@3\endcsname{-0.1834}
\expandafter\def\csname nzG@2@1@3\endcsname{0.2374}
\expandafter\def\csname nzB@2@1@3\endcsname{-0.2379}
\expandafter\def\csname neR@2@1@3\endcsname{-0.1442}
\expandafter\def\csname neG@2@1@3\endcsname{0.2885}
\expandafter\def\csname neB@2@1@3\endcsname{-0.2385}
\expandafter\def\csname nzR@2@1@4\endcsname{-0.1033}
\expandafter\def\csname nzG@2@1@4\endcsname{0.0755}
\expandafter\def\csname nzB@2@1@4\endcsname{0.0590}
\expandafter\def\csname neR@2@1@4\endcsname{-0.1055}
\expandafter\def\csname neG@2@1@4\endcsname{0.0151}
\expandafter\def\csname neB@2@1@4\endcsname{-0.0151}
\expandafter\def\csname nzR@2@1@5\endcsname{-0.1978}
\expandafter\def\csname nzG@2@1@5\endcsname{-0.1658}
\expandafter\def\csname nzB@2@1@5\endcsname{-0.1117}
\expandafter\def\csname neR@2@1@5\endcsname{-0.2445}
\expandafter\def\csname neG@2@1@5\endcsname{-0.1251}
\expandafter\def\csname neB@2@1@5\endcsname{-0.0876}
\expandafter\def\csname nzR@2@1@6\endcsname{-0.1545}
\expandafter\def\csname nzG@2@1@6\endcsname{0.1512}
\expandafter\def\csname nzB@2@1@6\endcsname{-0.0657}
\expandafter\def\csname neR@2@1@6\endcsname{-0.1696}
\expandafter\def\csname neG@2@1@6\endcsname{0.1340}
\expandafter\def\csname neB@2@1@6\endcsname{-0.0760}
\expandafter\def\csname nzR@2@1@7\endcsname{0.0803}
\expandafter\def\csname nzG@2@1@7\endcsname{0.2530}
\expandafter\def\csname nzB@2@1@7\endcsname{0.1024}
\expandafter\def\csname neR@2@1@7\endcsname{0.0996}
\expandafter\def\csname neG@2@1@7\endcsname{0.2372}
\expandafter\def\csname neB@2@1@7\endcsname{0.1445}
\expandafter\def\csname nzR@2@2@0\endcsname{-0.0818}
\expandafter\def\csname nzG@2@2@0\endcsname{0.1169}
\expandafter\def\csname nzB@2@2@0\endcsname{0.1919}
\expandafter\def\csname neR@2@2@0\endcsname{-0.0702}
\expandafter\def\csname neG@2@2@0\endcsname{0.1063}
\expandafter\def\csname neB@2@2@0\endcsname{0.1264}
\expandafter\def\csname nzR@2@2@1\endcsname{0.0012}
\expandafter\def\csname nzG@2@2@1\endcsname{0.0560}
\expandafter\def\csname nzB@2@2@1\endcsname{-0.0903}
\expandafter\def\csname neR@2@2@1\endcsname{-0.0664}
\expandafter\def\csname neG@2@2@1\endcsname{0.0935}
\expandafter\def\csname neB@2@2@1\endcsname{-0.0189}
\expandafter\def\csname nzR@2@2@2\endcsname{0.1122}
\expandafter\def\csname nzG@2@2@2\endcsname{-0.1635}
\expandafter\def\csname nzB@2@2@2\endcsname{0.1103}
\expandafter\def\csname neR@2@2@2\endcsname{0.0415}
\expandafter\def\csname neG@2@2@2\endcsname{-0.1445}
\expandafter\def\csname neB@2@2@2\endcsname{0.1647}
\expandafter\def\csname nzR@2@2@3\endcsname{0.2158}
\expandafter\def\csname nzG@2@2@3\endcsname{0.1277}
\expandafter\def\csname nzB@2@2@3\endcsname{-0.1747}
\expandafter\def\csname neR@2@2@3\endcsname{0.2606}
\expandafter\def\csname neG@2@2@3\endcsname{0.1765}
\expandafter\def\csname neB@2@2@3\endcsname{-0.1546}
\expandafter\def\csname nzR@2@2@4\endcsname{-0.1957}
\expandafter\def\csname nzG@2@2@4\endcsname{0.2540}
\expandafter\def\csname nzB@2@2@4\endcsname{0.0516}
\expandafter\def\csname neR@2@2@4\endcsname{-0.2032}
\expandafter\def\csname neG@2@2@4\endcsname{0.2061}
\expandafter\def\csname neB@2@2@4\endcsname{-0.0033}
\expandafter\def\csname nzR@2@2@5\endcsname{-0.2418}
\expandafter\def\csname nzG@2@2@5\endcsname{0.1976}
\expandafter\def\csname nzB@2@2@5\endcsname{0.0340}
\expandafter\def\csname neR@2@2@5\endcsname{-0.1793}
\expandafter\def\csname neG@2@2@5\endcsname{0.1446}
\expandafter\def\csname neB@2@2@5\endcsname{0.0740}
\expandafter\def\csname nzR@2@2@6\endcsname{-0.0536}
\expandafter\def\csname nzG@2@2@6\endcsname{0.1989}
\expandafter\def\csname nzB@2@2@6\endcsname{-0.0270}
\expandafter\def\csname neR@2@2@6\endcsname{-0.0076}
\expandafter\def\csname neG@2@2@6\endcsname{0.2423}
\expandafter\def\csname neB@2@2@6\endcsname{-0.0874}
\expandafter\def\csname nzR@2@2@7\endcsname{0.1491}
\expandafter\def\csname nzG@2@2@7\endcsname{-0.0533}
\expandafter\def\csname nzB@2@2@7\endcsname{-0.2522}
\expandafter\def\csname neR@2@2@7\endcsname{0.1374}
\expandafter\def\csname neG@2@2@7\endcsname{-0.0316}
\expandafter\def\csname neB@2@2@7\endcsname{-0.2603}
\expandafter\def\csname nzR@2@3@0\endcsname{-0.1356}
\expandafter\def\csname nzG@2@3@0\endcsname{0.1153}
\expandafter\def\csname nzB@2@3@0\endcsname{0.0306}
\expandafter\def\csname neR@2@3@0\endcsname{-0.1885}
\expandafter\def\csname neG@2@3@0\endcsname{0.0874}
\expandafter\def\csname neB@2@3@0\endcsname{0.0886}
\expandafter\def\csname nzR@2@3@1\endcsname{0.0591}
\expandafter\def\csname nzG@2@3@1\endcsname{0.0306}
\expandafter\def\csname nzB@2@3@1\endcsname{0.0801}
\expandafter\def\csname neR@2@3@1\endcsname{0.0231}
\expandafter\def\csname neG@2@3@1\endcsname{-0.0457}
\expandafter\def\csname neB@2@3@1\endcsname{0.0609}
\expandafter\def\csname nzR@2@3@2\endcsname{-0.0527}
\expandafter\def\csname nzG@2@3@2\endcsname{-0.1695}
\expandafter\def\csname nzB@2@3@2\endcsname{0.1164}
\expandafter\def\csname neR@2@3@2\endcsname{-0.1202}
\expandafter\def\csname neG@2@3@2\endcsname{-0.1908}
\expandafter\def\csname neB@2@3@2\endcsname{0.1357}
\expandafter\def\csname nzR@2@3@3\endcsname{0.0673}
\expandafter\def\csname nzG@2@3@3\endcsname{-0.2362}
\expandafter\def\csname nzB@2@3@3\endcsname{0.0156}
\expandafter\def\csname neR@2@3@3\endcsname{0.1243}
\expandafter\def\csname neG@2@3@3\endcsname{-0.2725}
\expandafter\def\csname neB@2@3@3\endcsname{-0.0488}
\expandafter\def\csname nzR@2@3@4\endcsname{-0.1042}
\expandafter\def\csname nzG@2@3@4\endcsname{0.2096}
\expandafter\def\csname nzB@2@3@4\endcsname{-0.1941}
\expandafter\def\csname neR@2@3@4\endcsname{-0.0556}
\expandafter\def\csname neG@2@3@4\endcsname{0.2389}
\expandafter\def\csname neB@2@3@4\endcsname{-0.1853}
\expandafter\def\csname nzR@2@3@5\endcsname{-0.0066}
\expandafter\def\csname nzG@2@3@5\endcsname{-0.0397}
\expandafter\def\csname nzB@2@3@5\endcsname{0.0026}
\expandafter\def\csname neR@2@3@5\endcsname{0.0683}
\expandafter\def\csname neG@2@3@5\endcsname{-0.0105}
\expandafter\def\csname neB@2@3@5\endcsname{-0.0095}
\expandafter\def\csname nzR@2@3@6\endcsname{0.1124}
\expandafter\def\csname nzG@2@3@6\endcsname{0.1848}
\expandafter\def\csname nzB@2@3@6\endcsname{-0.1071}
\expandafter\def\csname neR@2@3@6\endcsname{0.1282}
\expandafter\def\csname neG@2@3@6\endcsname{0.1678}
\expandafter\def\csname neB@2@3@6\endcsname{-0.1035}
\expandafter\def\csname nzR@2@3@7\endcsname{0.0632}
\expandafter\def\csname nzG@2@3@7\endcsname{0.1151}
\expandafter\def\csname nzB@2@3@7\endcsname{0.0537}
\expandafter\def\csname neR@2@3@7\endcsname{0.0415}
\expandafter\def\csname neG@2@3@7\endcsname{0.1811}
\expandafter\def\csname neB@2@3@7\endcsname{0.0415}
\expandafter\def\csname nzR@2@4@0\endcsname{0.2373}
\expandafter\def\csname nzG@2@4@0\endcsname{0.0274}
\expandafter\def\csname nzB@2@4@0\endcsname{0.1080}
\expandafter\def\csname neR@2@4@0\endcsname{0.1769}
\expandafter\def\csname neG@2@4@0\endcsname{0.0172}
\expandafter\def\csname neB@2@4@0\endcsname{0.1476}
\expandafter\def\csname nzR@2@4@1\endcsname{-0.2280}
\expandafter\def\csname nzG@2@4@1\endcsname{-0.1726}
\expandafter\def\csname nzB@2@4@1\endcsname{-0.1454}
\expandafter\def\csname neR@2@4@1\endcsname{-0.2462}
\expandafter\def\csname neG@2@4@1\endcsname{-0.2484}
\expandafter\def\csname neB@2@4@1\endcsname{-0.0815}
\expandafter\def\csname nzR@2@4@2\endcsname{-0.1273}
\expandafter\def\csname nzG@2@4@2\endcsname{-0.1898}
\expandafter\def\csname nzB@2@4@2\endcsname{0.0142}
\expandafter\def\csname neR@2@4@2\endcsname{-0.1786}
\expandafter\def\csname neG@2@4@2\endcsname{-0.2588}
\expandafter\def\csname neB@2@4@2\endcsname{0.0836}
\expandafter\def\csname nzR@2@4@3\endcsname{-0.1841}
\expandafter\def\csname nzG@2@4@3\endcsname{-0.0154}
\expandafter\def\csname nzB@2@4@3\endcsname{-0.1420}
\expandafter\def\csname neR@2@4@3\endcsname{-0.2183}
\expandafter\def\csname neG@2@4@3\endcsname{-0.0755}
\expandafter\def\csname neB@2@4@3\endcsname{-0.1079}
\expandafter\def\csname nzR@2@4@4\endcsname{0.1567}
\expandafter\def\csname nzG@2@4@4\endcsname{-0.2426}
\expandafter\def\csname nzB@2@4@4\endcsname{-0.2856}
\expandafter\def\csname neR@2@4@4\endcsname{0.1293}
\expandafter\def\csname neG@2@4@4\endcsname{-0.2430}
\expandafter\def\csname neB@2@4@4\endcsname{-0.3216}
\expandafter\def\csname nzR@2@4@5\endcsname{0.0266}
\expandafter\def\csname nzG@2@4@5\endcsname{-0.2362}
\expandafter\def\csname nzB@2@4@5\endcsname{0.0112}
\expandafter\def\csname neR@2@4@5\endcsname{0.0132}
\expandafter\def\csname neG@2@4@5\endcsname{-0.1943}
\expandafter\def\csname neB@2@4@5\endcsname{-0.0267}
\expandafter\def\csname nzR@2@4@6\endcsname{-0.2079}
\expandafter\def\csname nzG@2@4@6\endcsname{0.1206}
\expandafter\def\csname nzB@2@4@6\endcsname{0.1574}
\expandafter\def\csname neR@2@4@6\endcsname{-0.2522}
\expandafter\def\csname neG@2@4@6\endcsname{0.1446}
\expandafter\def\csname neB@2@4@6\endcsname{0.1782}
\expandafter\def\csname nzR@2@4@7\endcsname{-0.0808}
\expandafter\def\csname nzG@2@4@7\endcsname{-0.1600}
\expandafter\def\csname nzB@2@4@7\endcsname{-0.0767}
\expandafter\def\csname neR@2@4@7\endcsname{-0.0604}
\expandafter\def\csname neG@2@4@7\endcsname{-0.1055}
\expandafter\def\csname neB@2@4@7\endcsname{-0.0804}
\expandafter\def\csname nzR@2@5@0\endcsname{0.1875}
\expandafter\def\csname nzG@2@5@0\endcsname{-0.1585}
\expandafter\def\csname nzB@2@5@0\endcsname{0.2658}
\expandafter\def\csname neR@2@5@0\endcsname{0.1508}
\expandafter\def\csname neG@2@5@0\endcsname{-0.2290}
\expandafter\def\csname neB@2@5@0\endcsname{0.2270}
\expandafter\def\csname nzR@2@5@1\endcsname{0.0439}
\expandafter\def\csname nzG@2@5@1\endcsname{0.1714}
\expandafter\def\csname nzB@2@5@1\endcsname{-0.1991}
\expandafter\def\csname neR@2@5@1\endcsname{-0.0076}
\expandafter\def\csname neG@2@5@1\endcsname{0.0927}
\expandafter\def\csname neB@2@5@1\endcsname{-0.1811}
\expandafter\def\csname nzR@2@5@2\endcsname{-0.2315}
\expandafter\def\csname nzG@2@5@2\endcsname{-0.1789}
\expandafter\def\csname nzB@2@5@2\endcsname{-0.1884}
\expandafter\def\csname neR@2@5@2\endcsname{-0.2730}
\expandafter\def\csname neG@2@5@2\endcsname{-0.1697}
\expandafter\def\csname neB@2@5@2\endcsname{-0.1538}
\expandafter\def\csname nzR@2@5@3\endcsname{-0.2318}
\expandafter\def\csname nzG@2@5@3\endcsname{0.2233}
\expandafter\def\csname nzB@2@5@3\endcsname{0.1692}
\expandafter\def\csname neR@2@5@3\endcsname{-0.2080}
\expandafter\def\csname neG@2@5@3\endcsname{0.1830}
\expandafter\def\csname neB@2@5@3\endcsname{0.1319}
\expandafter\def\csname nzR@2@5@4\endcsname{-0.0515}
\expandafter\def\csname nzG@2@5@4\endcsname{0.0247}
\expandafter\def\csname nzB@2@5@4\endcsname{-0.2026}
\expandafter\def\csname neR@2@5@4\endcsname{-0.0976}
\expandafter\def\csname neG@2@5@4\endcsname{-0.0408}
\expandafter\def\csname neB@2@5@4\endcsname{-0.2376}
\expandafter\def\csname nzR@2@5@5\endcsname{0.2321}
\expandafter\def\csname nzG@2@5@5\endcsname{-0.2811}
\expandafter\def\csname nzB@2@5@5\endcsname{0.1346}
\expandafter\def\csname neR@2@5@5\endcsname{0.1637}
\expandafter\def\csname neG@2@5@5\endcsname{-0.3202}
\expandafter\def\csname neB@2@5@5\endcsname{0.1351}
\expandafter\def\csname nzR@2@5@6\endcsname{-0.0110}
\expandafter\def\csname nzG@2@5@6\endcsname{-0.2414}
\expandafter\def\csname nzB@2@5@6\endcsname{0.2723}
\expandafter\def\csname neR@2@5@6\endcsname{-0.0731}
\expandafter\def\csname neG@2@5@6\endcsname{-0.2924}
\expandafter\def\csname neB@2@5@6\endcsname{0.2773}
\expandafter\def\csname nzR@2@5@7\endcsname{-0.1344}
\expandafter\def\csname nzG@2@5@7\endcsname{-0.2242}
\expandafter\def\csname nzB@2@5@7\endcsname{0.0711}
\expandafter\def\csname neR@2@5@7\endcsname{-0.1542}
\expandafter\def\csname neG@2@5@7\endcsname{-0.2595}
\expandafter\def\csname neB@2@5@7\endcsname{0.1199}
\expandafter\def\csname nzR@2@6@0\endcsname{-0.2099}
\expandafter\def\csname nzG@2@6@0\endcsname{-0.1171}
\expandafter\def\csname nzB@2@6@0\endcsname{0.2772}
\expandafter\def\csname neR@2@6@0\endcsname{-0.1742}
\expandafter\def\csname neG@2@6@0\endcsname{-0.1786}
\expandafter\def\csname neB@2@6@0\endcsname{0.2633}
\expandafter\def\csname nzR@2@6@1\endcsname{0.2087}
\expandafter\def\csname nzG@2@6@1\endcsname{-0.0934}
\expandafter\def\csname nzB@2@6@1\endcsname{0.0097}
\expandafter\def\csname neR@2@6@1\endcsname{0.1328}
\expandafter\def\csname neG@2@6@1\endcsname{-0.1346}
\expandafter\def\csname neB@2@6@1\endcsname{-0.0145}
\expandafter\def\csname nzR@2@6@2\endcsname{0.1325}
\expandafter\def\csname nzG@2@6@2\endcsname{0.2237}
\expandafter\def\csname nzB@2@6@2\endcsname{-0.1649}
\expandafter\def\csname neR@2@6@2\endcsname{0.1055}
\expandafter\def\csname neG@2@6@2\endcsname{0.2328}
\expandafter\def\csname neB@2@6@2\endcsname{-0.1009}
\expandafter\def\csname nzR@2@6@3\endcsname{-0.0282}
\expandafter\def\csname nzG@2@6@3\endcsname{-0.0189}
\expandafter\def\csname nzB@2@6@3\endcsname{0.0555}
\expandafter\def\csname neR@2@6@3\endcsname{0.0159}
\expandafter\def\csname neG@2@6@3\endcsname{0.0406}
\expandafter\def\csname neB@2@6@3\endcsname{0.0620}
\expandafter\def\csname nzR@2@6@4\endcsname{-0.1859}
\expandafter\def\csname nzG@2@6@4\endcsname{-0.1328}
\expandafter\def\csname nzB@2@6@4\endcsname{-0.2110}
\expandafter\def\csname neR@2@6@4\endcsname{-0.1886}
\expandafter\def\csname neG@2@6@4\endcsname{-0.1802}
\expandafter\def\csname neB@2@6@4\endcsname{-0.1482}
\expandafter\def\csname nzR@2@6@5\endcsname{0.0998}
\expandafter\def\csname nzG@2@6@5\endcsname{-0.0567}
\expandafter\def\csname nzB@2@6@5\endcsname{0.1514}
\expandafter\def\csname neR@2@6@5\endcsname{0.1374}
\expandafter\def\csname neG@2@6@5\endcsname{-0.0793}
\expandafter\def\csname neB@2@6@5\endcsname{0.1044}
\expandafter\def\csname nzR@2@6@6\endcsname{0.1152}
\expandafter\def\csname nzG@2@6@6\endcsname{0.1577}
\expandafter\def\csname nzB@2@6@6\endcsname{-0.0588}
\expandafter\def\csname neR@2@6@6\endcsname{0.1060}
\expandafter\def\csname neG@2@6@6\endcsname{0.1247}
\expandafter\def\csname neB@2@6@6\endcsname{0.0033}
\expandafter\def\csname nzR@2@6@7\endcsname{-0.2390}
\expandafter\def\csname nzG@2@6@7\endcsname{0.2443}
\expandafter\def\csname nzB@2@6@7\endcsname{-0.0305}
\expandafter\def\csname neR@2@6@7\endcsname{-0.3023}
\expandafter\def\csname neG@2@6@7\endcsname{0.2915}
\expandafter\def\csname neB@2@6@7\endcsname{-0.0627}
\expandafter\def\csname nzR@2@7@0\endcsname{0.2556}
\expandafter\def\csname nzG@2@7@0\endcsname{0.0005}
\expandafter\def\csname nzB@2@7@0\endcsname{-0.0513}
\expandafter\def\csname neR@2@7@0\endcsname{0.1708}
\expandafter\def\csname neG@2@7@0\endcsname{-0.0513}
\expandafter\def\csname neB@2@7@0\endcsname{-0.0915}
\expandafter\def\csname nzR@2@7@1\endcsname{0.2097}
\expandafter\def\csname nzG@2@7@1\endcsname{0.1364}
\expandafter\def\csname nzB@2@7@1\endcsname{-0.1473}
\expandafter\def\csname neR@2@7@1\endcsname{0.1819}
\expandafter\def\csname neG@2@7@1\endcsname{0.1312}
\expandafter\def\csname neB@2@7@1\endcsname{-0.1783}
\expandafter\def\csname nzR@2@7@2\endcsname{-0.2009}
\expandafter\def\csname nzG@2@7@2\endcsname{0.0056}
\expandafter\def\csname nzB@2@7@2\endcsname{-0.0439}
\expandafter\def\csname neR@2@7@2\endcsname{-0.1362}
\expandafter\def\csname neG@2@7@2\endcsname{0.0592}
\expandafter\def\csname neB@2@7@2\endcsname{-0.0490}
\expandafter\def\csname nzR@2@7@3\endcsname{-0.1729}
\expandafter\def\csname nzG@2@7@3\endcsname{-0.0276}
\expandafter\def\csname nzB@2@7@3\endcsname{0.1784}
\expandafter\def\csname neR@2@7@3\endcsname{-0.1346}
\expandafter\def\csname neG@2@7@3\endcsname{0.0559}
\expandafter\def\csname neB@2@7@3\endcsname{0.1851}
\expandafter\def\csname nzR@2@7@4\endcsname{0.2298}
\expandafter\def\csname nzG@2@7@4\endcsname{0.0404}
\expandafter\def\csname nzB@2@7@4\endcsname{-0.0024}
\expandafter\def\csname neR@2@7@4\endcsname{0.2143}
\expandafter\def\csname neG@2@7@4\endcsname{0.1206}
\expandafter\def\csname neB@2@7@4\endcsname{-0.0426}
\expandafter\def\csname nzR@2@7@5\endcsname{0.0011}
\expandafter\def\csname nzG@2@7@5\endcsname{-0.1943}
\expandafter\def\csname nzB@2@7@5\endcsname{0.2025}
\expandafter\def\csname neR@2@7@5\endcsname{0.0407}
\expandafter\def\csname neG@2@7@5\endcsname{-0.1358}
\expandafter\def\csname neB@2@7@5\endcsname{0.2470}
\expandafter\def\csname nzR@2@7@6\endcsname{-0.1133}
\expandafter\def\csname nzG@2@7@6\endcsname{0.1289}
\expandafter\def\csname nzB@2@7@6\endcsname{-0.0904}
\expandafter\def\csname neR@2@7@6\endcsname{-0.1621}
\expandafter\def\csname neG@2@7@6\endcsname{0.0700}
\expandafter\def\csname neB@2@7@6\endcsname{-0.0762}
\expandafter\def\csname nzR@2@7@7\endcsname{0.2224}
\expandafter\def\csname nzG@2@7@7\endcsname{0.2763}
\expandafter\def\csname nzB@2@7@7\endcsname{0.2087}
\expandafter\def\csname neR@2@7@7\endcsname{0.1599}
\expandafter\def\csname neG@2@7@7\endcsname{0.3295}
\expandafter\def\csname neB@2@7@7\endcsname{0.1567}
\expandafter\def\csname nzR@3@0@0\endcsname{-0.2346}
\expandafter\def\csname nzG@3@0@0\endcsname{-0.0788}
\expandafter\def\csname nzB@3@0@0\endcsname{-0.3189}
\expandafter\def\csname neR@3@0@0\endcsname{-0.2254}
\expandafter\def\csname neG@3@0@0\endcsname{0.0175}
\expandafter\def\csname neB@3@0@0\endcsname{-0.3316}
\expandafter\def\csname nzR@3@0@1\endcsname{-0.0704}
\expandafter\def\csname nzG@3@0@1\endcsname{0.0198}
\expandafter\def\csname nzB@3@0@1\endcsname{-0.0255}
\expandafter\def\csname neR@3@0@1\endcsname{-0.1177}
\expandafter\def\csname neG@3@0@1\endcsname{0.0921}
\expandafter\def\csname neB@3@0@1\endcsname{-0.0427}
\expandafter\def\csname nzR@3@0@2\endcsname{-0.1447}
\expandafter\def\csname nzG@3@0@2\endcsname{-0.0407}
\expandafter\def\csname nzB@3@0@2\endcsname{0.0380}
\expandafter\def\csname neR@3@0@2\endcsname{-0.1365}
\expandafter\def\csname neG@3@0@2\endcsname{0.0261}
\expandafter\def\csname neB@3@0@2\endcsname{0.0297}
\expandafter\def\csname nzR@3@0@3\endcsname{-0.2227}
\expandafter\def\csname nzG@3@0@3\endcsname{-0.2255}
\expandafter\def\csname nzB@3@0@3\endcsname{-0.0275}
\expandafter\def\csname neR@3@0@3\endcsname{-0.1555}
\expandafter\def\csname neG@3@0@3\endcsname{-0.1715}
\expandafter\def\csname neB@3@0@3\endcsname{0.0116}
\expandafter\def\csname nzR@3@0@4\endcsname{0.0935}
\expandafter\def\csname nzG@3@0@4\endcsname{-0.2425}
\expandafter\def\csname nzB@3@0@4\endcsname{0.1357}
\expandafter\def\csname neR@3@0@4\endcsname{0.0488}
\expandafter\def\csname neG@3@0@4\endcsname{-0.2060}
\expandafter\def\csname neB@3@0@4\endcsname{0.1435}
\expandafter\def\csname nzR@3@0@5\endcsname{0.1096}
\expandafter\def\csname nzG@3@0@5\endcsname{0.0187}
\expandafter\def\csname nzB@3@0@5\endcsname{-0.0357}
\expandafter\def\csname neR@3@0@5\endcsname{0.0356}
\expandafter\def\csname neG@3@0@5\endcsname{0.0245}
\expandafter\def\csname neB@3@0@5\endcsname{-0.0223}
\expandafter\def\csname nzR@3@0@6\endcsname{-0.0532}
\expandafter\def\csname nzG@3@0@6\endcsname{0.1240}
\expandafter\def\csname nzB@3@0@6\endcsname{0.2065}
\expandafter\def\csname neR@3@0@6\endcsname{-0.1321}
\expandafter\def\csname neG@3@0@6\endcsname{0.1740}
\expandafter\def\csname neB@3@0@6\endcsname{0.1513}
\expandafter\def\csname nzR@3@0@7\endcsname{-0.2474}
\expandafter\def\csname nzG@3@0@7\endcsname{-0.0748}
\expandafter\def\csname nzB@3@0@7\endcsname{0.4112}
\expandafter\def\csname neR@3@0@7\endcsname{-0.2713}
\expandafter\def\csname neG@3@0@7\endcsname{0.0253}
\expandafter\def\csname neB@3@0@7\endcsname{0.3224}
\expandafter\def\csname nzR@3@1@0\endcsname{0.0575}
\expandafter\def\csname nzG@3@1@0\endcsname{0.0200}
\expandafter\def\csname nzB@3@1@0\endcsname{-0.2374}
\expandafter\def\csname neR@3@1@0\endcsname{0.0059}
\expandafter\def\csname neG@3@1@0\endcsname{0.0568}
\expandafter\def\csname neB@3@1@0\endcsname{-0.2831}
\expandafter\def\csname nzR@3@1@1\endcsname{0.0920}
\expandafter\def\csname nzG@3@1@1\endcsname{0.1038}
\expandafter\def\csname nzB@3@1@1\endcsname{-0.1261}
\expandafter\def\csname neR@3@1@1\endcsname{0.0206}
\expandafter\def\csname neG@3@1@1\endcsname{0.1139}
\expandafter\def\csname neB@3@1@1\endcsname{-0.1259}
\expandafter\def\csname nzR@3@1@2\endcsname{0.0541}
\expandafter\def\csname nzG@3@1@2\endcsname{-0.0336}
\expandafter\def\csname nzB@3@1@2\endcsname{0.0014}
\expandafter\def\csname neR@3@1@2\endcsname{0.0424}
\expandafter\def\csname neG@3@1@2\endcsname{-0.0012}
\expandafter\def\csname neB@3@1@2\endcsname{0.0392}
\expandafter\def\csname nzR@3@1@3\endcsname{-0.1248}
\expandafter\def\csname nzG@3@1@3\endcsname{0.1178}
\expandafter\def\csname nzB@3@1@3\endcsname{-0.1506}
\expandafter\def\csname neR@3@1@3\endcsname{-0.0757}
\expandafter\def\csname neG@3@1@3\endcsname{0.1478}
\expandafter\def\csname neB@3@1@3\endcsname{-0.1341}
\expandafter\def\csname nzR@3@1@4\endcsname{-0.1533}
\expandafter\def\csname nzG@3@1@4\endcsname{0.0748}
\expandafter\def\csname nzB@3@1@4\endcsname{0.0060}
\expandafter\def\csname neR@3@1@4\endcsname{-0.1515}
\expandafter\def\csname neG@3@1@4\endcsname{0.0464}
\expandafter\def\csname neB@3@1@4\endcsname{-0.0373}
\expandafter\def\csname nzR@3@1@5\endcsname{-0.2091}
\expandafter\def\csname nzG@3@1@5\endcsname{0.0064}
\expandafter\def\csname nzB@3@1@5\endcsname{-0.0689}
\expandafter\def\csname neR@3@1@5\endcsname{-0.2290}
\expandafter\def\csname neG@3@1@5\endcsname{0.0074}
\expandafter\def\csname neB@3@1@5\endcsname{-0.0617}
\expandafter\def\csname nzR@3@1@6\endcsname{-0.1493}
\expandafter\def\csname nzG@3@1@6\endcsname{0.1837}
\expandafter\def\csname nzB@3@1@6\endcsname{-0.0150}
\expandafter\def\csname neR@3@1@6\endcsname{-0.1614}
\expandafter\def\csname neG@3@1@6\endcsname{0.1869}
\expandafter\def\csname neB@3@1@6\endcsname{-0.0282}
\expandafter\def\csname nzR@3@1@7\endcsname{0.0017}
\expandafter\def\csname nzG@3@1@7\endcsname{0.2104}
\expandafter\def\csname nzB@3@1@7\endcsname{0.0923}
\expandafter\def\csname neR@3@1@7\endcsname{0.0040}
\expandafter\def\csname neG@3@1@7\endcsname{0.2122}
\expandafter\def\csname neB@3@1@7\endcsname{0.0936}
\expandafter\def\csname nzR@3@2@0\endcsname{-0.0530}
\expandafter\def\csname nzG@3@2@0\endcsname{0.1430}
\expandafter\def\csname nzB@3@2@0\endcsname{0.0821}
\expandafter\def\csname neR@3@2@0\endcsname{-0.0849}
\expandafter\def\csname neG@3@2@0\endcsname{0.1285}
\expandafter\def\csname neB@3@2@0\endcsname{0.0415}
\expandafter\def\csname nzR@3@2@1\endcsname{0.0436}
\expandafter\def\csname nzG@3@2@1\endcsname{0.0596}
\expandafter\def\csname nzB@3@2@1\endcsname{-0.0159}
\expandafter\def\csname neR@3@2@1\endcsname{-0.0388}
\expandafter\def\csname neG@3@2@1\endcsname{0.0620}
\expandafter\def\csname neB@3@2@1\endcsname{0.0244}
\expandafter\def\csname nzR@3@2@2\endcsname{0.1354}
\expandafter\def\csname nzG@3@2@2\endcsname{-0.1210}
\expandafter\def\csname nzB@3@2@2\endcsname{0.0455}
\expandafter\def\csname neR@3@2@2\endcsname{0.0691}
\expandafter\def\csname neG@3@2@2\endcsname{-0.0932}
\expandafter\def\csname neB@3@2@2\endcsname{0.1010}
\expandafter\def\csname nzR@3@2@3\endcsname{0.0998}
\expandafter\def\csname nzG@3@2@3\endcsname{0.1044}
\expandafter\def\csname nzB@3@2@3\endcsname{-0.1236}
\expandafter\def\csname neR@3@2@3\endcsname{0.1253}
\expandafter\def\csname neG@3@2@3\endcsname{0.1240}
\expandafter\def\csname neB@3@2@3\endcsname{-0.1165}
\expandafter\def\csname nzR@3@2@4\endcsname{-0.1983}
\expandafter\def\csname nzG@3@2@4\endcsname{0.2830}
\expandafter\def\csname nzB@3@2@4\endcsname{-0.0391}
\expandafter\def\csname neR@3@2@4\endcsname{-0.1553}
\expandafter\def\csname neG@3@2@4\endcsname{0.2335}
\expandafter\def\csname neB@3@2@4\endcsname{-0.0755}
\expandafter\def\csname nzR@3@2@5\endcsname{-0.2700}
\expandafter\def\csname nzG@3@2@5\endcsname{0.2183}
\expandafter\def\csname nzB@3@2@5\endcsname{-0.0094}
\expandafter\def\csname neR@3@2@5\endcsname{-0.1992}
\expandafter\def\csname neG@3@2@5\endcsname{0.1796}
\expandafter\def\csname neB@3@2@5\endcsname{-0.0066}
\expandafter\def\csname nzR@3@2@6\endcsname{-0.0673}
\expandafter\def\csname nzG@3@2@6\endcsname{0.2378}
\expandafter\def\csname nzB@3@2@6\endcsname{-0.0927}
\expandafter\def\csname neR@3@2@6\endcsname{-0.0221}
\expandafter\def\csname neG@3@2@6\endcsname{0.2449}
\expandafter\def\csname neB@3@2@6\endcsname{-0.1182}
\expandafter\def\csname nzR@3@2@7\endcsname{0.1435}
\expandafter\def\csname nzG@3@2@7\endcsname{0.1204}
\expandafter\def\csname nzB@3@2@7\endcsname{-0.1897}
\expandafter\def\csname neR@3@2@7\endcsname{0.1360}
\expandafter\def\csname neG@3@2@7\endcsname{0.1483}
\expandafter\def\csname neB@3@2@7\endcsname{-0.1924}
\expandafter\def\csname nzR@3@3@0\endcsname{-0.0773}
\expandafter\def\csname nzG@3@3@0\endcsname{0.1298}
\expandafter\def\csname nzB@3@3@0\endcsname{0.0991}
\expandafter\def\csname neR@3@3@0\endcsname{-0.1374}
\expandafter\def\csname neG@3@3@0\endcsname{0.0759}
\expandafter\def\csname neB@3@3@0\endcsname{0.1369}
\expandafter\def\csname nzR@3@3@1\endcsname{-0.0346}
\expandafter\def\csname nzG@3@3@1\endcsname{-0.0281}
\expandafter\def\csname nzB@3@3@1\endcsname{0.0621}
\expandafter\def\csname neR@3@3@1\endcsname{-0.1026}
\expandafter\def\csname neG@3@3@1\endcsname{-0.0944}
\expandafter\def\csname neB@3@3@1\endcsname{0.0911}
\expandafter\def\csname nzR@3@3@2\endcsname{-0.0228}
\expandafter\def\csname nzG@3@3@2\endcsname{-0.2266}
\expandafter\def\csname nzB@3@3@2\endcsname{0.0916}
\expandafter\def\csname neR@3@3@2\endcsname{-0.0869}
\expandafter\def\csname neG@3@3@2\endcsname{-0.2594}
\expandafter\def\csname neB@3@3@2\endcsname{0.1158}
\expandafter\def\csname nzR@3@3@3\endcsname{0.0200}
\expandafter\def\csname nzG@3@3@3\endcsname{-0.1504}
\expandafter\def\csname nzB@3@3@3\endcsname{-0.0736}
\expandafter\def\csname neR@3@3@3\endcsname{0.0445}
\expandafter\def\csname neG@3@3@3\endcsname{-0.1700}
\expandafter\def\csname neB@3@3@3\endcsname{-0.0926}
\expandafter\def\csname nzR@3@3@4\endcsname{-0.0752}
\expandafter\def\csname nzG@3@3@4\endcsname{0.0923}
\expandafter\def\csname nzB@3@3@4\endcsname{-0.1890}
\expandafter\def\csname neR@3@3@4\endcsname{-0.0179}
\expandafter\def\csname neG@3@3@4\endcsname{0.0945}
\expandafter\def\csname neB@3@3@4\endcsname{-0.1999}
\expandafter\def\csname nzR@3@3@5\endcsname{-0.0634}
\expandafter\def\csname nzG@3@3@5\endcsname{0.0660}
\expandafter\def\csname nzB@3@3@5\endcsname{-0.0545}
\expandafter\def\csname neR@3@3@5\endcsname{0.0076}
\expandafter\def\csname neG@3@3@5\endcsname{0.0818}
\expandafter\def\csname neB@3@3@5\endcsname{-0.0630}
\expandafter\def\csname nzR@3@3@6\endcsname{0.0270}
\expandafter\def\csname nzG@3@3@6\endcsname{0.1892}
\expandafter\def\csname nzB@3@3@6\endcsname{-0.0515}
\expandafter\def\csname neR@3@3@6\endcsname{0.0496}
\expandafter\def\csname neG@3@3@6\endcsname{0.2003}
\expandafter\def\csname neB@3@3@6\endcsname{-0.0585}
\expandafter\def\csname nzR@3@3@7\endcsname{0.0776}
\expandafter\def\csname nzG@3@3@7\endcsname{0.1131}
\expandafter\def\csname nzB@3@3@7\endcsname{-0.0539}
\expandafter\def\csname neR@3@3@7\endcsname{0.0616}
\expandafter\def\csname neG@3@3@7\endcsname{0.1751}
\expandafter\def\csname neB@3@3@7\endcsname{-0.0651}
\expandafter\def\csname nzR@3@4@0\endcsname{0.1593}
\expandafter\def\csname nzG@3@4@0\endcsname{-0.0145}
\expandafter\def\csname nzB@3@4@0\endcsname{0.1112}
\expandafter\def\csname neR@3@4@0\endcsname{0.0732}
\expandafter\def\csname neG@3@4@0\endcsname{-0.0715}
\expandafter\def\csname neB@3@4@0\endcsname{0.1541}
\expandafter\def\csname nzR@3@4@1\endcsname{-0.1294}
\expandafter\def\csname nzG@3@4@1\endcsname{-0.1325}
\expandafter\def\csname nzB@3@4@1\endcsname{-0.0880}
\expandafter\def\csname neR@3@4@1\endcsname{-0.1799}
\expandafter\def\csname neG@3@4@1\endcsname{-0.2222}
\expandafter\def\csname neB@3@4@1\endcsname{-0.0186}
\expandafter\def\csname nzR@3@4@2\endcsname{-0.2308}
\expandafter\def\csname nzG@3@4@2\endcsname{-0.2228}
\expandafter\def\csname nzB@3@4@2\endcsname{-0.0562}
\expandafter\def\csname neR@3@4@2\endcsname{-0.2789}
\expandafter\def\csname neG@3@4@2\endcsname{-0.2935}
\expandafter\def\csname neB@3@4@2\endcsname{0.0131}
\expandafter\def\csname nzR@3@4@3\endcsname{-0.1768}
\expandafter\def\csname nzG@3@4@3\endcsname{-0.1037}
\expandafter\def\csname nzB@3@4@3\endcsname{-0.1385}
\expandafter\def\csname neR@3@4@3\endcsname{-0.1933}
\expandafter\def\csname neG@3@4@3\endcsname{-0.1654}
\expandafter\def\csname neB@3@4@3\endcsname{-0.1202}
\expandafter\def\csname nzR@3@4@4\endcsname{0.0505}
\expandafter\def\csname nzG@3@4@4\endcsname{-0.1864}
\expandafter\def\csname nzB@3@4@4\endcsname{-0.2785}
\expandafter\def\csname neR@3@4@4\endcsname{0.0270}
\expandafter\def\csname neG@3@4@4\endcsname{-0.1897}
\expandafter\def\csname neB@3@4@4\endcsname{-0.2970}
\expandafter\def\csname nzR@3@4@5\endcsname{0.0518}
\expandafter\def\csname nzG@3@4@5\endcsname{-0.2371}
\expandafter\def\csname nzB@3@4@5\endcsname{-0.0041}
\expandafter\def\csname neR@3@4@5\endcsname{0.0252}
\expandafter\def\csname neG@3@4@5\endcsname{-0.1990}
\expandafter\def\csname neB@3@4@5\endcsname{-0.0338}
\expandafter\def\csname nzR@3@4@6\endcsname{-0.1221}
\expandafter\def\csname nzG@3@4@6\endcsname{-0.0323}
\expandafter\def\csname nzB@3@4@6\endcsname{0.1454}
\expandafter\def\csname neR@3@4@6\endcsname{-0.1527}
\expandafter\def\csname neG@3@4@6\endcsname{-0.0094}
\expandafter\def\csname neB@3@4@6\endcsname{0.1454}
\expandafter\def\csname nzR@3@4@7\endcsname{-0.1276}
\expandafter\def\csname nzG@3@4@7\endcsname{-0.1301}
\expandafter\def\csname nzB@3@4@7\endcsname{0.0204}
\expandafter\def\csname neR@3@4@7\endcsname{-0.1270}
\expandafter\def\csname neG@3@4@7\endcsname{-0.0733}
\expandafter\def\csname neB@3@4@7\endcsname{0.0307}
\expandafter\def\csname nzR@3@5@0\endcsname{0.1692}
\expandafter\def\csname nzG@3@5@0\endcsname{-0.1272}
\expandafter\def\csname nzB@3@5@0\endcsname{0.2512}
\expandafter\def\csname neR@3@5@0\endcsname{0.1079}
\expandafter\def\csname neG@3@5@0\endcsname{-0.2159}
\expandafter\def\csname neB@3@5@0\endcsname{0.2188}
\expandafter\def\csname nzR@3@5@1\endcsname{0.0166}
\expandafter\def\csname nzG@3@5@1\endcsname{-0.0066}
\expandafter\def\csname nzB@3@5@1\endcsname{-0.1396}
\expandafter\def\csname neR@3@5@1\endcsname{-0.0546}
\expandafter\def\csname neG@3@5@1\endcsname{-0.0954}
\expandafter\def\csname neB@3@5@1\endcsname{-0.1035}
\expandafter\def\csname nzR@3@5@2\endcsname{-0.2051}
\expandafter\def\csname nzG@3@5@2\endcsname{-0.0510}
\expandafter\def\csname nzB@3@5@2\endcsname{-0.1761}
\expandafter\def\csname neR@3@5@2\endcsname{-0.2431}
\expandafter\def\csname neG@3@5@2\endcsname{-0.0813}
\expandafter\def\csname neB@3@5@2\endcsname{-0.1177}
\expandafter\def\csname nzR@3@5@3\endcsname{-0.2563}
\expandafter\def\csname nzG@3@5@3\endcsname{0.0955}
\expandafter\def\csname nzB@3@5@3\endcsname{-0.0187}
\expandafter\def\csname neR@3@5@3\endcsname{-0.2462}
\expandafter\def\csname neG@3@5@3\endcsname{0.0480}
\expandafter\def\csname neB@3@5@3\endcsname{-0.0246}
\expandafter\def\csname nzR@3@5@4\endcsname{-0.0447}
\expandafter\def\csname nzG@3@5@4\endcsname{-0.0871}
\expandafter\def\csname nzB@3@5@4\endcsname{-0.1702}
\expandafter\def\csname neR@3@5@4\endcsname{-0.0831}
\expandafter\def\csname neG@3@5@4\endcsname{-0.1465}
\expandafter\def\csname neB@3@5@4\endcsname{-0.1868}
\expandafter\def\csname nzR@3@5@5\endcsname{0.1639}
\expandafter\def\csname nzG@3@5@5\endcsname{-0.2958}
\expandafter\def\csname nzB@3@5@5\endcsname{0.1122}
\expandafter\def\csname neR@3@5@5\endcsname{0.0901}
\expandafter\def\csname neG@3@5@5\endcsname{-0.3243}
\expandafter\def\csname neB@3@5@5\endcsname{0.0903}
\expandafter\def\csname nzR@3@5@6\endcsname{-0.0148}
\expandafter\def\csname nzG@3@5@6\endcsname{-0.2199}
\expandafter\def\csname nzB@3@5@6\endcsname{0.2508}
\expandafter\def\csname neR@3@5@6\endcsname{-0.0804}
\expandafter\def\csname neG@3@5@6\endcsname{-0.2479}
\expandafter\def\csname neB@3@5@6\endcsname{0.2568}
\expandafter\def\csname nzR@3@5@7\endcsname{-0.1992}
\expandafter\def\csname nzG@3@5@7\endcsname{-0.2032}
\expandafter\def\csname nzB@3@5@7\endcsname{0.1052}
\expandafter\def\csname neR@3@5@7\endcsname{-0.2297}
\expandafter\def\csname neG@3@5@7\endcsname{-0.2082}
\expandafter\def\csname neB@3@5@7\endcsname{0.1356}
\expandafter\def\csname nzR@3@6@0\endcsname{-0.0007}
\expandafter\def\csname nzG@3@6@0\endcsname{-0.1371}
\expandafter\def\csname nzB@3@6@0\endcsname{0.2599}
\expandafter\def\csname neR@3@6@0\endcsname{-0.0223}
\expandafter\def\csname neG@3@6@0\endcsname{-0.2191}
\expandafter\def\csname neB@3@6@0\endcsname{0.2091}
\expandafter\def\csname nzR@3@6@1\endcsname{0.1791}
\expandafter\def\csname nzG@3@6@1\endcsname{0.0011}
\expandafter\def\csname nzB@3@6@1\endcsname{-0.0445}
\expandafter\def\csname neR@3@6@1\endcsname{0.1008}
\expandafter\def\csname neG@3@6@1\endcsname{-0.0533}
\expandafter\def\csname neB@3@6@1\endcsname{-0.0525}
\expandafter\def\csname nzR@3@6@2\endcsname{0.0360}
\expandafter\def\csname nzG@3@6@2\endcsname{0.1264}
\expandafter\def\csname nzB@3@6@2\endcsname{-0.1446}
\expandafter\def\csname neR@3@6@2\endcsname{0.0134}
\expandafter\def\csname neG@3@6@2\endcsname{0.1377}
\expandafter\def\csname neB@3@6@2\endcsname{-0.0956}
\expandafter\def\csname nzR@3@6@3\endcsname{-0.1247}
\expandafter\def\csname nzG@3@6@3\endcsname{0.0399}
\expandafter\def\csname nzB@3@6@3\endcsname{0.0043}
\expandafter\def\csname neR@3@6@3\endcsname{-0.0853}
\expandafter\def\csname neG@3@6@3\endcsname{0.0797}
\expandafter\def\csname neB@3@6@3\endcsname{0.0240}
\expandafter\def\csname nzR@3@6@4\endcsname{-0.0902}
\expandafter\def\csname nzG@3@6@4\endcsname{-0.1082}
\expandafter\def\csname nzB@3@6@4\endcsname{-0.1017}
\expandafter\def\csname neR@3@6@4\endcsname{-0.0805}
\expandafter\def\csname neG@3@6@4\endcsname{-0.1204}
\expandafter\def\csname neB@3@6@4\endcsname{-0.0766}
\expandafter\def\csname nzR@3@6@5\endcsname{0.1047}
\expandafter\def\csname nzG@3@6@5\endcsname{-0.1293}
\expandafter\def\csname nzB@3@6@5\endcsname{0.1114}
\expandafter\def\csname neR@3@6@5\endcsname{0.1053}
\expandafter\def\csname neG@3@6@5\endcsname{-0.1544}
\expandafter\def\csname neB@3@6@5\endcsname{0.1035}
\expandafter\def\csname nzR@3@6@6\endcsname{0.0507}
\expandafter\def\csname nzG@3@6@6\endcsname{0.0928}
\expandafter\def\csname nzB@3@6@6\endcsname{0.0571}
\expandafter\def\csname neR@3@6@6\endcsname{0.0119}
\expandafter\def\csname neG@3@6@6\endcsname{0.0515}
\expandafter\def\csname neB@3@6@6\endcsname{0.0847}
\expandafter\def\csname nzR@3@6@7\endcsname{-0.1592}
\expandafter\def\csname nzG@3@6@7\endcsname{0.2417}
\expandafter\def\csname nzB@3@6@7\endcsname{0.0367}
\expandafter\def\csname neR@3@6@7\endcsname{-0.2263}
\expandafter\def\csname neG@3@6@7\endcsname{0.2515}
\expandafter\def\csname neB@3@6@7\endcsname{0.0259}
\expandafter\def\csname nzR@3@7@0\endcsname{0.2717}
\expandafter\def\csname nzG@3@7@0\endcsname{-0.0022}
\expandafter\def\csname nzB@3@7@0\endcsname{-0.0110}
\expandafter\def\csname neR@3@7@0\endcsname{0.1723}
\expandafter\def\csname neG@3@7@0\endcsname{-0.0726}
\expandafter\def\csname neB@3@7@0\endcsname{-0.0624}
\expandafter\def\csname nzR@3@7@1\endcsname{0.2305}
\expandafter\def\csname nzG@3@7@1\endcsname{0.0984}
\expandafter\def\csname nzB@3@7@1\endcsname{-0.1344}
\expandafter\def\csname neR@3@7@1\endcsname{0.1746}
\expandafter\def\csname neG@3@7@1\endcsname{0.0765}
\expandafter\def\csname neB@3@7@1\endcsname{-0.1623}
\expandafter\def\csname nzR@3@7@2\endcsname{-0.1089}
\expandafter\def\csname nzG@3@7@2\endcsname{0.0788}
\expandafter\def\csname nzB@3@7@2\endcsname{-0.0633}
\expandafter\def\csname neR@3@7@2\endcsname{-0.0582}
\expandafter\def\csname neG@3@7@2\endcsname{0.1364}
\expandafter\def\csname neB@3@7@2\endcsname{-0.0548}
\expandafter\def\csname nzR@3@7@3\endcsname{-0.1450}
\expandafter\def\csname nzG@3@7@3\endcsname{-0.0100}
\expandafter\def\csname nzB@3@7@3\endcsname{0.1194}
\expandafter\def\csname neR@3@7@3\endcsname{-0.0876}
\expandafter\def\csname neG@3@7@3\endcsname{0.0963}
\expandafter\def\csname neB@3@7@3\endcsname{0.1178}
\expandafter\def\csname nzR@3@7@4\endcsname{0.1069}
\expandafter\def\csname nzG@3@7@4\endcsname{-0.0607}
\expandafter\def\csname nzB@3@7@4\endcsname{0.0536}
\expandafter\def\csname neR@3@7@4\endcsname{0.1115}
\expandafter\def\csname neG@3@7@4\endcsname{0.0285}
\expandafter\def\csname neB@3@7@4\endcsname{0.0430}
\expandafter\def\csname nzR@3@7@5\endcsname{0.0476}
\expandafter\def\csname nzG@3@7@5\endcsname{-0.1300}
\expandafter\def\csname nzB@3@7@5\endcsname{0.1577}
\expandafter\def\csname neR@3@7@5\endcsname{0.0693}
\expandafter\def\csname neG@3@7@5\endcsname{-0.0839}
\expandafter\def\csname neB@3@7@5\endcsname{0.1744}
\expandafter\def\csname nzR@3@7@6\endcsname{-0.0244}
\expandafter\def\csname nzG@3@7@6\endcsname{0.1760}
\expandafter\def\csname nzB@3@7@6\endcsname{0.0144}
\expandafter\def\csname neR@3@7@6\endcsname{-0.0680}
\expandafter\def\csname neG@3@7@6\endcsname{0.1381}
\expandafter\def\csname neB@3@7@6\endcsname{0.0331}
\expandafter\def\csname nzR@3@7@7\endcsname{0.1244}
\expandafter\def\csname nzG@3@7@7\endcsname{0.4047}
\expandafter\def\csname nzB@3@7@7\endcsname{0.1610}
\expandafter\def\csname neR@3@7@7\endcsname{0.0260}
\expandafter\def\csname neG@3@7@7\endcsname{0.4238}
\expandafter\def\csname neB@3@7@7\endcsname{0.1027}
\expandafter\def\csname nzR@4@0@0\endcsname{-0.1656}
\expandafter\def\csname nzG@4@0@0\endcsname{-0.0379}
\expandafter\def\csname nzB@4@0@0\endcsname{-0.3031}
\expandafter\def\csname neR@4@0@0\endcsname{-0.1784}
\expandafter\def\csname neG@4@0@0\endcsname{0.0557}
\expandafter\def\csname neB@4@0@0\endcsname{-0.3090}
\expandafter\def\csname nzR@4@0@1\endcsname{-0.0839}
\expandafter\def\csname nzG@4@0@1\endcsname{0.0098}
\expandafter\def\csname nzB@4@0@1\endcsname{-0.0969}
\expandafter\def\csname neR@4@0@1\endcsname{-0.1206}
\expandafter\def\csname neG@4@0@1\endcsname{0.0840}
\expandafter\def\csname neB@4@0@1\endcsname{-0.1047}
\expandafter\def\csname nzR@4@0@2\endcsname{-0.1326}
\expandafter\def\csname nzG@4@0@2\endcsname{-0.0553}
\expandafter\def\csname nzB@4@0@2\endcsname{-0.0078}
\expandafter\def\csname neR@4@0@2\endcsname{-0.1205}
\expandafter\def\csname neG@4@0@2\endcsname{0.0159}
\expandafter\def\csname neB@4@0@2\endcsname{0.0010}
\expandafter\def\csname nzR@4@0@3\endcsname{-0.1678}
\expandafter\def\csname nzG@4@0@3\endcsname{-0.1597}
\expandafter\def\csname nzB@4@0@3\endcsname{-0.0087}
\expandafter\def\csname neR@4@0@3\endcsname{-0.1205}
\expandafter\def\csname neG@4@0@3\endcsname{-0.0962}
\expandafter\def\csname neB@4@0@3\endcsname{0.0150}
\expandafter\def\csname nzR@4@0@4\endcsname{-0.0229}
\expandafter\def\csname nzG@4@0@4\endcsname{-0.1557}
\expandafter\def\csname nzB@4@0@4\endcsname{0.0492}
\expandafter\def\csname neR@4@0@4\endcsname{-0.0440}
\expandafter\def\csname neG@4@0@4\endcsname{-0.1195}
\expandafter\def\csname neB@4@0@4\endcsname{0.0543}
\expandafter\def\csname nzR@4@0@5\endcsname{0.0105}
\expandafter\def\csname nzG@4@0@5\endcsname{0.0078}
\expandafter\def\csname nzB@4@0@5\endcsname{0.0343}
\expandafter\def\csname neR@4@0@5\endcsname{-0.0577}
\expandafter\def\csname neG@4@0@5\endcsname{0.0249}
\expandafter\def\csname neB@4@0@5\endcsname{0.0281}
\expandafter\def\csname nzR@4@0@6\endcsname{-0.0958}
\expandafter\def\csname nzG@4@0@6\endcsname{0.1163}
\expandafter\def\csname nzB@4@0@6\endcsname{0.1779}
\expandafter\def\csname neR@4@0@6\endcsname{-0.1556}
\expandafter\def\csname neG@4@0@6\endcsname{0.1553}
\expandafter\def\csname neB@4@0@6\endcsname{0.1240}
\expandafter\def\csname nzR@4@0@7\endcsname{-0.2017}
\expandafter\def\csname nzG@4@0@7\endcsname{0.0437}
\expandafter\def\csname nzB@4@0@7\endcsname{0.3659}
\expandafter\def\csname neR@4@0@7\endcsname{-0.2254}
\expandafter\def\csname neG@4@0@7\endcsname{0.1215}
\expandafter\def\csname neB@4@0@7\endcsname{0.2714}
\expandafter\def\csname nzR@4@1@0\endcsname{0.0021}
\expandafter\def\csname nzG@4@1@0\endcsname{0.0542}
\expandafter\def\csname nzB@4@1@0\endcsname{-0.2042}
\expandafter\def\csname neR@4@1@0\endcsname{-0.0533}
\expandafter\def\csname neG@4@1@0\endcsname{0.0903}
\expandafter\def\csname neB@4@1@0\endcsname{-0.2261}
\expandafter\def\csname nzR@4@1@1\endcsname{0.0501}
\expandafter\def\csname nzG@4@1@1\endcsname{0.0541}
\expandafter\def\csname nzB@4@1@1\endcsname{-0.1025}
\expandafter\def\csname neR@4@1@1\endcsname{-0.0187}
\expandafter\def\csname neG@4@1@1\endcsname{0.0830}
\expandafter\def\csname neB@4@1@1\endcsname{-0.0914}
\expandafter\def\csname nzR@4@1@2\endcsname{0.0181}
\expandafter\def\csname nzG@4@1@2\endcsname{-0.0043}
\expandafter\def\csname nzB@4@1@2\endcsname{-0.0444}
\expandafter\def\csname neR@4@1@2\endcsname{-0.0031}
\expandafter\def\csname neG@4@1@2\endcsname{0.0362}
\expandafter\def\csname neB@4@1@2\endcsname{-0.0102}
\expandafter\def\csname nzR@4@1@3\endcsname{-0.0933}
\expandafter\def\csname nzG@4@1@3\endcsname{0.0462}
\expandafter\def\csname nzB@4@1@3\endcsname{-0.0844}
\expandafter\def\csname neR@4@1@3\endcsname{-0.0584}
\expandafter\def\csname neG@4@1@3\endcsname{0.0705}
\expandafter\def\csname neB@4@1@3\endcsname{-0.0663}
\expandafter\def\csname nzR@4@1@4\endcsname{-0.1629}
\expandafter\def\csname nzG@4@1@4\endcsname{0.0802}
\expandafter\def\csname nzB@4@1@4\endcsname{-0.0340}
\expandafter\def\csname neR@4@1@4\endcsname{-0.1423}
\expandafter\def\csname neG@4@1@4\endcsname{0.0660}
\expandafter\def\csname neB@4@1@4\endcsname{-0.0491}
\expandafter\def\csname nzR@4@1@5\endcsname{-0.1925}
\expandafter\def\csname nzG@4@1@5\endcsname{0.1084}
\expandafter\def\csname nzB@4@1@5\endcsname{-0.0303}
\expandafter\def\csname neR@4@1@5\endcsname{-0.1912}
\expandafter\def\csname neG@4@1@5\endcsname{0.0944}
\expandafter\def\csname neB@4@1@5\endcsname{-0.0379}
\expandafter\def\csname nzR@4@1@6\endcsname{-0.1426}
\expandafter\def\csname nzG@4@1@6\endcsname{0.1944}
\expandafter\def\csname nzB@4@1@6\endcsname{0.0088}
\expandafter\def\csname neR@4@1@6\endcsname{-0.1462}
\expandafter\def\csname neG@4@1@6\endcsname{0.1961}
\expandafter\def\csname neB@4@1@6\endcsname{-0.0099}
\expandafter\def\csname nzR@4@1@7\endcsname{-0.0495}
\expandafter\def\csname nzG@4@1@7\endcsname{0.1976}
\expandafter\def\csname nzB@4@1@7\endcsname{0.0845}
\expandafter\def\csname neR@4@1@7\endcsname{-0.0541}
\expandafter\def\csname neG@4@1@7\endcsname{0.2123}
\expandafter\def\csname neB@4@1@7\endcsname{0.0588}
\expandafter\def\csname nzR@4@2@0\endcsname{-0.0202}
\expandafter\def\csname nzG@4@2@0\endcsname{0.1281}
\expandafter\def\csname nzB@4@2@0\endcsname{0.0130}
\expandafter\def\csname neR@4@2@0\endcsname{-0.0817}
\expandafter\def\csname neG@4@2@0\endcsname{0.1062}
\expandafter\def\csname neB@4@2@0\endcsname{-0.0018}
\expandafter\def\csname nzR@4@2@1\endcsname{0.0472}
\expandafter\def\csname nzG@4@2@1\endcsname{0.0315}
\expandafter\def\csname nzB@4@2@1\endcsname{0.0010}
\expandafter\def\csname neR@4@2@1\endcsname{-0.0375}
\expandafter\def\csname neG@4@2@1\endcsname{0.0233}
\expandafter\def\csname neB@4@2@1\endcsname{0.0298}
\expandafter\def\csname nzR@4@2@2\endcsname{0.0884}
\expandafter\def\csname nzG@4@2@2\endcsname{-0.0678}
\expandafter\def\csname nzB@4@2@2\endcsname{-0.0034}
\expandafter\def\csname neR@4@2@2\endcsname{0.0320}
\expandafter\def\csname neG@4@2@2\endcsname{-0.0511}
\expandafter\def\csname neB@4@2@2\endcsname{0.0384}
\expandafter\def\csname nzR@4@2@3\endcsname{0.0084}
\expandafter\def\csname nzG@4@2@3\endcsname{0.0642}
\expandafter\def\csname nzB@4@2@3\endcsname{-0.0941}
\expandafter\def\csname neR@4@2@3\endcsname{0.0255}
\expandafter\def\csname neG@4@2@3\endcsname{0.0657}
\expandafter\def\csname neB@4@2@3\endcsname{-0.0837}
\expandafter\def\csname nzR@4@2@4\endcsname{-0.1746}
\expandafter\def\csname nzG@4@2@4\endcsname{0.2102}
\expandafter\def\csname nzB@4@2@4\endcsname{-0.0881}
\expandafter\def\csname neR@4@2@4\endcsname{-0.1160}
\expandafter\def\csname neG@4@2@4\endcsname{0.1715}
\expandafter\def\csname neB@4@2@4\endcsname{-0.1055}
\expandafter\def\csname nzR@4@2@5\endcsname{-0.2288}
\expandafter\def\csname nzG@4@2@5\endcsname{0.2234}
\expandafter\def\csname nzB@4@2@5\endcsname{-0.0574}
\expandafter\def\csname neR@4@2@5\endcsname{-0.1609}
\expandafter\def\csname neG@4@2@5\endcsname{0.1891}
\expandafter\def\csname neB@4@2@5\endcsname{-0.0664}
\expandafter\def\csname nzR@4@2@6\endcsname{-0.0887}
\expandafter\def\csname nzG@4@2@6\endcsname{0.2437}
\expandafter\def\csname nzB@4@2@6\endcsname{-0.0879}
\expandafter\def\csname neR@4@2@6\endcsname{-0.0491}
\expandafter\def\csname neG@4@2@6\endcsname{0.2361}
\expandafter\def\csname neB@4@2@6\endcsname{-0.0972}
\expandafter\def\csname nzR@4@2@7\endcsname{0.0765}
\expandafter\def\csname nzG@4@2@7\endcsname{0.2066}
\expandafter\def\csname nzB@4@2@7\endcsname{-0.1244}
\expandafter\def\csname neR@4@2@7\endcsname{0.0749}
\expandafter\def\csname neG@4@2@7\endcsname{0.2233}
\expandafter\def\csname neB@4@2@7\endcsname{-0.1262}
\expandafter\def\csname nzR@4@3@0\endcsname{-0.0375}
\expandafter\def\csname nzG@4@3@0\endcsname{0.0875}
\expandafter\def\csname nzB@4@3@0\endcsname{0.1014}
\expandafter\def\csname neR@4@3@0\endcsname{-0.1107}
\expandafter\def\csname neG@4@3@0\endcsname{0.0226}
\expandafter\def\csname neB@4@3@0\endcsname{0.1292}
\expandafter\def\csname nzR@4@3@1\endcsname{-0.0464}
\expandafter\def\csname nzG@4@3@1\endcsname{-0.0606}
\expandafter\def\csname nzB@4@3@1\endcsname{0.0501}
\expandafter\def\csname neR@4@3@1\endcsname{-0.1235}
\expandafter\def\csname neG@4@3@1\endcsname{-0.1170}
\expandafter\def\csname neB@4@3@1\endcsname{0.0929}
\expandafter\def\csname nzR@4@3@2\endcsname{-0.0359}
\expandafter\def\csname nzG@4@3@2\endcsname{-0.1914}
\expandafter\def\csname nzB@4@3@2\endcsname{0.0182}
\expandafter\def\csname neR@4@3@2\endcsname{-0.0911}
\expandafter\def\csname neG@4@3@2\endcsname{-0.2202}
\expandafter\def\csname neB@4@3@2\endcsname{0.0549}
\expandafter\def\csname nzR@4@3@3\endcsname{-0.0308}
\expandafter\def\csname nzG@4@3@3\endcsname{-0.1035}
\expandafter\def\csname nzB@4@3@3\endcsname{-0.1081}
\expandafter\def\csname neR@4@3@3\endcsname{-0.0236}
\expandafter\def\csname neG@4@3@3\endcsname{-0.1215}
\expandafter\def\csname neB@4@3@3\endcsname{-0.1003}
\expandafter\def\csname nzR@4@3@4\endcsname{-0.0794}
\expandafter\def\csname nzG@4@3@4\endcsname{0.0415}
\expandafter\def\csname nzB@4@3@4\endcsname{-0.1718}
\expandafter\def\csname neR@4@3@4\endcsname{-0.0300}
\expandafter\def\csname neG@4@3@4\endcsname{0.0298}
\expandafter\def\csname neB@4@3@4\endcsname{-0.1789}
\expandafter\def\csname nzR@4@3@5\endcsname{-0.0900}
\expandafter\def\csname nzG@4@3@5\endcsname{0.0899}
\expandafter\def\csname nzB@4@3@5\endcsname{-0.0793}
\expandafter\def\csname neR@4@3@5\endcsname{-0.0341}
\expandafter\def\csname neG@4@3@5\endcsname{0.0922}
\expandafter\def\csname neB@4@3@5\endcsname{-0.0892}
\expandafter\def\csname nzR@4@3@6\endcsname{-0.0305}
\expandafter\def\csname nzG@4@3@6\endcsname{0.1473}
\expandafter\def\csname nzB@4@3@6\endcsname{-0.0400}
\expandafter\def\csname neR@4@3@6\endcsname{-0.0048}
\expandafter\def\csname neG@4@3@6\endcsname{0.1628}
\expandafter\def\csname neB@4@3@6\endcsname{-0.0488}
\expandafter\def\csname nzR@4@3@7\endcsname{0.0388}
\expandafter\def\csname nzG@4@3@7\endcsname{0.1135}
\expandafter\def\csname nzB@4@3@7\endcsname{-0.0723}
\expandafter\def\csname neR@4@3@7\endcsname{0.0322}
\expandafter\def\csname neG@4@3@7\endcsname{0.1610}
\expandafter\def\csname neB@4@3@7\endcsname{-0.0762}
\expandafter\def\csname nzR@4@4@0\endcsname{0.0832}
\expandafter\def\csname nzG@4@4@0\endcsname{-0.0367}
\expandafter\def\csname nzB@4@4@0\endcsname{0.1089}
\expandafter\def\csname neR@4@4@0\endcsname{-0.0095}
\expandafter\def\csname neG@4@4@0\endcsname{-0.1188}
\expandafter\def\csname neB@4@4@0\endcsname{0.1440}
\expandafter\def\csname nzR@4@4@1\endcsname{-0.0898}
\expandafter\def\csname nzG@4@4@1\endcsname{-0.1232}
\expandafter\def\csname nzB@4@4@1\endcsname{-0.0388}
\expandafter\def\csname neR@4@4@1\endcsname{-0.1579}
\expandafter\def\csname neG@4@4@1\endcsname{-0.2079}
\expandafter\def\csname neB@4@4@1\endcsname{0.0251}
\expandafter\def\csname nzR@4@4@2\endcsname{-0.2058}
\expandafter\def\csname nzG@4@4@2\endcsname{-0.1877}
\expandafter\def\csname nzB@4@4@2\endcsname{-0.0834}
\expandafter\def\csname neR@4@4@2\endcsname{-0.2446}
\expandafter\def\csname neG@4@4@2\endcsname{-0.2509}
\expandafter\def\csname neB@4@4@2\endcsname{-0.0221}
\expandafter\def\csname nzR@4@4@3\endcsname{-0.1662}
\expandafter\def\csname nzG@4@4@3\endcsname{-0.1349}
\expandafter\def\csname nzB@4@4@3\endcsname{-0.1538}
\expandafter\def\csname neR@4@4@3\endcsname{-0.1737}
\expandafter\def\csname neG@4@4@3\endcsname{-0.1808}
\expandafter\def\csname neB@4@4@3\endcsname{-0.1309}
\expandafter\def\csname nzR@4@4@4\endcsname{-0.0224}
\expandafter\def\csname nzG@4@4@4\endcsname{-0.1537}
\expandafter\def\csname nzB@4@4@4\endcsname{-0.2074}
\expandafter\def\csname neR@4@4@4\endcsname{-0.0328}
\expandafter\def\csname neG@4@4@4\endcsname{-0.1623}
\expandafter\def\csname neB@4@4@4\endcsname{-0.2128}
\expandafter\def\csname nzR@4@4@5\endcsname{0.0235}
\expandafter\def\csname nzG@4@4@5\endcsname{-0.1811}
\expandafter\def\csname nzB@4@4@5\endcsname{-0.0232}
\expandafter\def\csname neR@4@4@5\endcsname{0.0005}
\expandafter\def\csname neG@4@4@5\endcsname{-0.1572}
\expandafter\def\csname neB@4@4@5\endcsname{-0.0433}
\expandafter\def\csname nzR@4@4@6\endcsname{-0.0631}
\expandafter\def\csname nzG@4@4@6\endcsname{-0.0961}
\expandafter\def\csname nzB@4@4@6\endcsname{0.1111}
\expandafter\def\csname neR@4@4@6\endcsname{-0.0862}
\expandafter\def\csname neG@4@4@6\endcsname{-0.0684}
\expandafter\def\csname neB@4@4@6\endcsname{0.1004}
\expandafter\def\csname nzR@4@4@7\endcsname{-0.1187}
\expandafter\def\csname nzG@4@4@7\endcsname{-0.1027}
\expandafter\def\csname nzB@4@4@7\endcsname{0.0651}
\expandafter\def\csname neR@4@4@7\endcsname{-0.1289}
\expandafter\def\csname neG@4@4@7\endcsname{-0.0526}
\expandafter\def\csname neB@4@4@7\endcsname{0.0704}
\expandafter\def\csname nzR@4@5@0\endcsname{0.1326}
\expandafter\def\csname nzG@4@5@0\endcsname{-0.1142}
\expandafter\def\csname nzB@4@5@0\endcsname{0.1815}
\expandafter\def\csname neR@4@5@0\endcsname{0.0508}
\expandafter\def\csname neG@4@5@0\endcsname{-0.2069}
\expandafter\def\csname neB@4@5@0\endcsname{0.1632}
\expandafter\def\csname nzR@4@5@1\endcsname{-0.0001}
\expandafter\def\csname nzG@4@5@1\endcsname{-0.0645}
\expandafter\def\csname nzB@4@5@1\endcsname{-0.0795}
\expandafter\def\csname neR@4@5@1\endcsname{-0.0740}
\expandafter\def\csname neG@4@5@1\endcsname{-0.1475}
\expandafter\def\csname neB@4@5@1\endcsname{-0.0375}
\expandafter\def\csname nzR@4@5@2\endcsname{-0.1762}
\expandafter\def\csname nzG@4@5@2\endcsname{-0.0344}
\expandafter\def\csname nzB@4@5@2\endcsname{-0.1518}
\expandafter\def\csname neR@4@5@2\endcsname{-0.2076}
\expandafter\def\csname neG@4@5@2\endcsname{-0.0832}
\expandafter\def\csname neB@4@5@2\endcsname{-0.0903}
\expandafter\def\csname nzR@4@5@3\endcsname{-0.2170}
\expandafter\def\csname nzG@4@5@3\endcsname{-0.0098}
\expandafter\def\csname nzB@4@5@3\endcsname{-0.1179}
\expandafter\def\csname neR@4@5@3\endcsname{-0.2131}
\expandafter\def\csname neG@4@5@3\endcsname{-0.0504}
\expandafter\def\csname neB@4@5@3\endcsname{-0.0948}
\expandafter\def\csname nzR@4@5@4\endcsname{-0.0490}
\expandafter\def\csname nzG@4@5@4\endcsname{-0.1416}
\expandafter\def\csname nzB@4@5@4\endcsname{-0.1243}
\expandafter\def\csname neR@4@5@4\endcsname{-0.0744}
\expandafter\def\csname neG@4@5@4\endcsname{-0.1735}
\expandafter\def\csname neB@4@5@4\endcsname{-0.1278}
\expandafter\def\csname nzR@4@5@5\endcsname{0.0841}
\expandafter\def\csname nzG@4@5@5\endcsname{-0.2545}
\expandafter\def\csname nzB@4@5@5\endcsname{0.0690}
\expandafter\def\csname neR@4@5@5\endcsname{0.0235}
\expandafter\def\csname neG@4@5@5\endcsname{-0.2655}
\expandafter\def\csname neB@4@5@5\endcsname{0.0504}
\expandafter\def\csname nzR@4@5@6\endcsname{-0.0226}
\expandafter\def\csname nzG@4@5@6\endcsname{-0.1865}
\expandafter\def\csname nzB@4@5@6\endcsname{0.1945}
\expandafter\def\csname neR@4@5@6\endcsname{-0.0806}
\expandafter\def\csname neG@4@5@6\endcsname{-0.1925}
\expandafter\def\csname neB@4@5@6\endcsname{0.1885}
\expandafter\def\csname nzR@4@5@7\endcsname{-0.1792}
\expandafter\def\csname nzG@4@5@7\endcsname{-0.1359}
\expandafter\def\csname nzB@4@5@7\endcsname{0.1328}
\expandafter\def\csname neR@4@5@7\endcsname{-0.2144}
\expandafter\def\csname neG@4@5@7\endcsname{-0.1247}
\expandafter\def\csname neB@4@5@7\endcsname{0.1444}
\expandafter\def\csname nzR@4@6@0\endcsname{0.1297}
\expandafter\def\csname nzG@4@6@0\endcsname{-0.1020}
\expandafter\def\csname nzB@4@6@0\endcsname{0.1779}
\expandafter\def\csname neR@4@6@0\endcsname{0.0600}
\expandafter\def\csname neG@4@6@0\endcsname{-0.1880}
\expandafter\def\csname neB@4@6@0\endcsname{0.1244}
\expandafter\def\csname nzR@4@6@1\endcsname{0.1297}
\expandafter\def\csname nzG@4@6@1\endcsname{0.0110}
\expandafter\def\csname nzB@4@6@1\endcsname{-0.0571}
\expandafter\def\csname neR@4@6@1\endcsname{0.0544}
\expandafter\def\csname neG@4@6@1\endcsname{-0.0466}
\expandafter\def\csname neB@4@6@1\endcsname{-0.0535}
\expandafter\def\csname nzR@4@6@2\endcsname{-0.0322}
\expandafter\def\csname nzG@4@6@2\endcsname{0.0786}
\expandafter\def\csname nzB@4@6@2\endcsname{-0.1213}
\expandafter\def\csname neR@4@6@2\endcsname{-0.0490}
\expandafter\def\csname neG@4@6@2\endcsname{0.0761}
\expandafter\def\csname neB@4@6@2\endcsname{-0.0806}
\expandafter\def\csname nzR@4@6@3\endcsname{-0.1456}
\expandafter\def\csname nzG@4@6@3\endcsname{0.0304}
\expandafter\def\csname nzB@4@6@3\endcsname{-0.0465}
\expandafter\def\csname neR@4@6@3\endcsname{-0.1146}
\expandafter\def\csname neG@4@6@3\endcsname{0.0535}
\expandafter\def\csname neB@4@6@3\endcsname{-0.0226}
\expandafter\def\csname nzR@4@6@4\endcsname{-0.0488}
\expandafter\def\csname nzG@4@6@4\endcsname{-0.1055}
\expandafter\def\csname nzB@4@6@4\endcsname{-0.0295}
\expandafter\def\csname neR@4@6@4\endcsname{-0.0408}
\expandafter\def\csname neG@4@6@4\endcsname{-0.0983}
\expandafter\def\csname neB@4@6@4\endcsname{-0.0204}
\expandafter\def\csname nzR@4@6@5\endcsname{0.0780}
\expandafter\def\csname nzG@4@6@5\endcsname{-0.1428}
\expandafter\def\csname nzB@4@6@5\endcsname{0.0907}
\expandafter\def\csname neR@4@6@5\endcsname{0.0530}
\expandafter\def\csname neG@4@6@5\endcsname{-0.1560}
\expandafter\def\csname neB@4@6@5\endcsname{0.0891}
\expandafter\def\csname nzR@4@6@6\endcsname{0.0238}
\expandafter\def\csname nzG@4@6@6\endcsname{0.0231}
\expandafter\def\csname nzB@4@6@6\endcsname{0.1254}
\expandafter\def\csname neR@4@6@6\endcsname{-0.0283}
\expandafter\def\csname neG@4@6@6\endcsname{-0.0087}
\expandafter\def\csname neB@4@6@6\endcsname{0.1302}
\expandafter\def\csname nzR@4@6@7\endcsname{-0.1062}
\expandafter\def\csname nzG@4@6@7\endcsname{0.1793}
\expandafter\def\csname nzB@4@6@7\endcsname{0.1019}
\expandafter\def\csname neR@4@6@7\endcsname{-0.1734}
\expandafter\def\csname neG@4@6@7\endcsname{0.1646}
\expandafter\def\csname neB@4@6@7\endcsname{0.0935}
\expandafter\def\csname nzR@4@7@0\endcsname{0.2664}
\expandafter\def\csname nzG@4@7@0\endcsname{-0.0110}
\expandafter\def\csname nzB@4@7@0\endcsname{0.0196}
\expandafter\def\csname neR@4@7@0\endcsname{0.1597}
\expandafter\def\csname neG@4@7@0\endcsname{-0.0864}
\expandafter\def\csname neB@4@7@0\endcsname{-0.0350}
\expandafter\def\csname nzR@4@7@1\endcsname{0.1978}
\expandafter\def\csname nzG@4@7@1\endcsname{0.0758}
\expandafter\def\csname nzB@4@7@1\endcsname{-0.1014}
\expandafter\def\csname neR@4@7@1\endcsname{0.1316}
\expandafter\def\csname neG@4@7@1\endcsname{0.0450}
\expandafter\def\csname neB@4@7@1\endcsname{-0.1182}
\expandafter\def\csname nzR@4@7@2\endcsname{-0.0371}
\expandafter\def\csname nzG@4@7@2\endcsname{0.0895}
\expandafter\def\csname nzB@4@7@2\endcsname{-0.0698}
\expandafter\def\csname neR@4@7@2\endcsname{-0.0137}
\expandafter\def\csname neG@4@7@2\endcsname{0.1326}
\expandafter\def\csname neB@4@7@2\endcsname{-0.0567}
\expandafter\def\csname nzR@4@7@3\endcsname{-0.1137}
\expandafter\def\csname nzG@4@7@3\endcsname{0.0078}
\expandafter\def\csname nzB@4@7@3\endcsname{0.0485}
\expandafter\def\csname neR@4@7@3\endcsname{-0.0597}
\expandafter\def\csname neG@4@7@3\endcsname{0.0988}
\expandafter\def\csname neB@4@7@3\endcsname{0.0530}
\expandafter\def\csname nzR@4@7@4\endcsname{0.0236}
\expandafter\def\csname nzG@4@7@4\endcsname{-0.0912}
\expandafter\def\csname nzB@4@7@4\endcsname{0.0754}
\expandafter\def\csname neR@4@7@4\endcsname{0.0436}
\expandafter\def\csname neG@4@7@4\endcsname{-0.0105}
\expandafter\def\csname neB@4@7@4\endcsname{0.0730}
\expandafter\def\csname nzR@4@7@5\endcsname{0.0622}
\expandafter\def\csname nzG@4@7@5\endcsname{-0.0802}
\expandafter\def\csname nzB@4@7@5\endcsname{0.1204}
\expandafter\def\csname neR@4@7@5\endcsname{0.0609}
\expandafter\def\csname neG@4@7@5\endcsname{-0.0497}
\expandafter\def\csname neB@4@7@5\endcsname{0.1263}
\expandafter\def\csname nzR@4@7@6\endcsname{0.0240}
\expandafter\def\csname nzG@4@7@6\endcsname{0.1649}
\expandafter\def\csname nzB@4@7@6\endcsname{0.0929}
\expandafter\def\csname neR@4@7@6\endcsname{-0.0238}
\expandafter\def\csname neG@4@7@6\endcsname{0.1376}
\expandafter\def\csname neB@4@7@6\endcsname{0.0964}
\expandafter\def\csname nzR@4@7@7\endcsname{0.0484}
\expandafter\def\csname nzG@4@7@7\endcsname{0.3976}
\expandafter\def\csname nzB@4@7@7\endcsname{0.1351}
\expandafter\def\csname neR@4@7@7\endcsname{-0.0497}
\expandafter\def\csname neG@4@7@7\endcsname{0.3786}
\expandafter\def\csname neB@4@7@7\endcsname{0.0921}
\expandafter\def\csname nhR@0@0\endcsname{-0.0364}
\expandafter\def\csname nhG@0@0\endcsname{-0.1059}
\expandafter\def\csname nhB@0@0\endcsname{0.0505}
\expandafter\def\csname nhR@0@1\endcsname{0.0904}
\expandafter\def\csname nhG@0@1\endcsname{-0.1009}
\expandafter\def\csname nhB@0@1\endcsname{-0.0298}
\expandafter\def\csname nhR@0@2\endcsname{0.0466}
\expandafter\def\csname nhG@0@2\endcsname{-0.1109}
\expandafter\def\csname nhB@0@2\endcsname{0.0841}
\expandafter\def\csname nhR@0@3\endcsname{-0.1324}
\expandafter\def\csname nhG@0@3\endcsname{0.0330}
\expandafter\def\csname nhB@0@3\endcsname{-0.1067}
\expandafter\def\csname nhR@0@4\endcsname{0.1108}
\expandafter\def\csname nhG@0@4\endcsname{-0.0682}
\expandafter\def\csname nhB@0@4\endcsname{-0.0593}
\expandafter\def\csname nhR@0@5\endcsname{0.0006}
\expandafter\def\csname nhG@0@5\endcsname{0.0618}
\expandafter\def\csname nhB@0@5\endcsname{-0.0259}
\expandafter\def\csname nhR@0@6\endcsname{0.1490}
\expandafter\def\csname nhG@0@6\endcsname{-0.1254}
\expandafter\def\csname nhB@0@6\endcsname{0.0562}
\expandafter\def\csname nhR@0@7\endcsname{0.0548}
\expandafter\def\csname nhG@0@7\endcsname{-0.1269}
\expandafter\def\csname nhB@0@7\endcsname{0.0862}
\expandafter\def\csname nhR@1@0\endcsname{0.0692}
\expandafter\def\csname nhG@1@0\endcsname{-0.0454}
\expandafter\def\csname nhB@1@0\endcsname{0.0995}
\expandafter\def\csname nhR@1@1\endcsname{0.0821}
\expandafter\def\csname nhG@1@1\endcsname{0.0418}
\expandafter\def\csname nhB@1@1\endcsname{0.0912}
\expandafter\def\csname nhR@1@2\endcsname{-0.0614}
\expandafter\def\csname nhG@1@2\endcsname{0.0294}
\expandafter\def\csname nhB@1@2\endcsname{-0.1161}
\expandafter\def\csname nhR@1@3\endcsname{-0.0323}
\expandafter\def\csname nhG@1@3\endcsname{-0.1379}
\expandafter\def\csname nhB@1@3\endcsname{0.0400}
\expandafter\def\csname nhR@1@4\endcsname{-0.0030}
\expandafter\def\csname nhG@1@4\endcsname{0.1315}
\expandafter\def\csname nhB@1@4\endcsname{0.1405}
\expandafter\def\csname nhR@1@5\endcsname{0.1217}
\expandafter\def\csname nhG@1@5\endcsname{-0.0885}
\expandafter\def\csname nhB@1@5\endcsname{-0.0396}
\expandafter\def\csname nhR@1@6\endcsname{0.0380}
\expandafter\def\csname nhG@1@6\endcsname{0.0439}
\expandafter\def\csname nhB@1@6\endcsname{0.0253}
\expandafter\def\csname nhR@1@7\endcsname{-0.0548}
\expandafter\def\csname nhG@1@7\endcsname{0.0108}
\expandafter\def\csname nhB@1@7\endcsname{-0.1048}
\expandafter\def\csname nhR@2@0\endcsname{-0.0382}
\expandafter\def\csname nhG@2@0\endcsname{0.0154}
\expandafter\def\csname nhB@2@0\endcsname{0.1138}
\expandafter\def\csname nhR@2@1\endcsname{0.1179}
\expandafter\def\csname nhG@2@1\endcsname{-0.0992}
\expandafter\def\csname nhB@2@1\endcsname{-0.1465}
\expandafter\def\csname nhR@2@2\endcsname{0.1234}
\expandafter\def\csname nhG@2@2\endcsname{-0.0063}
\expandafter\def\csname nhB@2@2\endcsname{-0.0965}
\expandafter\def\csname nhR@2@3\endcsname{-0.1190}
\expandafter\def\csname nhG@2@3\endcsname{-0.1173}
\expandafter\def\csname nhB@2@3\endcsname{-0.0352}
\expandafter\def\csname nhR@2@4\endcsname{0.0646}
\expandafter\def\csname nhG@2@4\endcsname{0.0654}
\expandafter\def\csname nhB@2@4\endcsname{0.1055}
\expandafter\def\csname nhR@2@5\endcsname{-0.0936}
\expandafter\def\csname nhG@2@5\endcsname{0.0989}
\expandafter\def\csname nhB@2@5\endcsname{-0.1047}
\expandafter\def\csname nhR@2@6\endcsname{-0.0790}
\expandafter\def\csname nhG@2@6\endcsname{-0.1248}
\expandafter\def\csname nhB@2@6\endcsname{0.1296}
\expandafter\def\csname nhR@2@7\endcsname{0.0131}
\expandafter\def\csname nhG@2@7\endcsname{-0.0145}
\expandafter\def\csname nhB@2@7\endcsname{0.0407}
\expandafter\def\csname nhR@3@0\endcsname{0.1080}
\expandafter\def\csname nhG@3@0\endcsname{0.0265}
\expandafter\def\csname nhB@3@0\endcsname{-0.1195}
\expandafter\def\csname nhR@3@1\endcsname{0.0300}
\expandafter\def\csname nhG@3@1\endcsname{0.1410}
\expandafter\def\csname nhB@3@1\endcsname{0.0652}
\expandafter\def\csname nhR@3@2\endcsname{0.1322}
\expandafter\def\csname nhG@3@2\endcsname{0.0312}
\expandafter\def\csname nhB@3@2\endcsname{-0.0497}
\expandafter\def\csname nhR@3@3\endcsname{-0.1287}
\expandafter\def\csname nhG@3@3\endcsname{0.1081}
\expandafter\def\csname nhB@3@3\endcsname{0.1432}
\expandafter\def\csname nhR@3@4\endcsname{-0.0701}
\expandafter\def\csname nhG@3@4\endcsname{-0.0974}
\expandafter\def\csname nhB@3@4\endcsname{-0.0250}
\expandafter\def\csname nhR@3@5\endcsname{-0.1355}
\expandafter\def\csname nhG@3@5\endcsname{-0.0523}
\expandafter\def\csname nhB@3@5\endcsname{0.0272}
\expandafter\def\csname nhR@3@6\endcsname{-0.0396}
\expandafter\def\csname nhG@3@6\endcsname{0.0438}
\expandafter\def\csname nhB@3@6\endcsname{-0.0028}
\expandafter\def\csname nhR@3@7\endcsname{0.0394}
\expandafter\def\csname nhG@3@7\endcsname{-0.1312}
\expandafter\def\csname nhB@3@7\endcsname{0.0105}
\expandafter\def\csname nhR@4@0\endcsname{0.0696}
\expandafter\def\csname nhG@4@0\endcsname{-0.0111}
\expandafter\def\csname nhB@4@0\endcsname{-0.0727}
\expandafter\def\csname nhR@4@1\endcsname{0.0398}
\expandafter\def\csname nhG@4@1\endcsname{0.1410}
\expandafter\def\csname nhB@4@1\endcsname{-0.0984}
\expandafter\def\csname nhR@4@2\endcsname{0.0889}
\expandafter\def\csname nhG@4@2\endcsname{0.1388}
\expandafter\def\csname nhB@4@2\endcsname{-0.1229}
\expandafter\def\csname nhR@4@3\endcsname{0.0939}
\expandafter\def\csname nhG@4@3\endcsname{0.0961}
\expandafter\def\csname nhB@4@3\endcsname{-0.0664}
\expandafter\def\csname nhR@4@4\endcsname{0.0267}
\expandafter\def\csname nhG@4@4\endcsname{0.0251}
\expandafter\def\csname nhB@4@4\endcsname{0.0998}
\expandafter\def\csname nhR@4@5\endcsname{0.0146}
\expandafter\def\csname nhG@4@5\endcsname{-0.0571}
\expandafter\def\csname nhB@4@5\endcsname{0.0730}
\expandafter\def\csname nhR@4@6\endcsname{0.1116}
\expandafter\def\csname nhG@4@6\endcsname{-0.0659}
\expandafter\def\csname nhB@4@6\endcsname{-0.0677}
\expandafter\def\csname nhR@4@7\endcsname{-0.0453}
\expandafter\def\csname nhG@4@7\endcsname{-0.0767}
\expandafter\def\csname nhB@4@7\endcsname{0.0294}
\expandafter\def\csname nhR@5@0\endcsname{0.0400}
\expandafter\def\csname nhG@5@0\endcsname{0.1314}
\expandafter\def\csname nhB@5@0\endcsname{0.0514}
\expandafter\def\csname nhR@5@1\endcsname{0.0725}
\expandafter\def\csname nhG@5@1\endcsname{0.1082}
\expandafter\def\csname nhB@5@1\endcsname{-0.0075}
\expandafter\def\csname nhR@5@2\endcsname{0.0973}
\expandafter\def\csname nhG@5@2\endcsname{-0.0202}
\expandafter\def\csname nhB@5@2\endcsname{-0.0302}
\expandafter\def\csname nhR@5@3\endcsname{-0.0375}
\expandafter\def\csname nhG@5@3\endcsname{0.0469}
\expandafter\def\csname nhB@5@3\endcsname{0.0567}
\expandafter\def\csname nhR@5@4\endcsname{0.0895}
\expandafter\def\csname nhG@5@4\endcsname{0.1049}
\expandafter\def\csname nhB@5@4\endcsname{0.0953}
\expandafter\def\csname nhR@5@5\endcsname{0.0985}
\expandafter\def\csname nhG@5@5\endcsname{0.0952}
\expandafter\def\csname nhB@5@5\endcsname{-0.0381}
\expandafter\def\csname nhR@5@6\endcsname{0.1073}
\expandafter\def\csname nhG@5@6\endcsname{0.1205}
\expandafter\def\csname nhB@5@6\endcsname{-0.0165}
\expandafter\def\csname nhR@5@7\endcsname{0.0349}
\expandafter\def\csname nhG@5@7\endcsname{0.0982}
\expandafter\def\csname nhB@5@7\endcsname{-0.1057}
\expandafter\def\csname nhR@6@0\endcsname{-0.0678}
\expandafter\def\csname nhG@6@0\endcsname{0.1030}
\expandafter\def\csname nhB@6@0\endcsname{-0.0208}
\expandafter\def\csname nhR@6@1\endcsname{0.1218}
\expandafter\def\csname nhG@6@1\endcsname{0.0799}
\expandafter\def\csname nhB@6@1\endcsname{0.0534}
\expandafter\def\csname nhR@6@2\endcsname{0.0355}
\expandafter\def\csname nhG@6@2\endcsname{-0.0337}
\expandafter\def\csname nhB@6@2\endcsname{-0.1132}
\expandafter\def\csname nhR@6@3\endcsname{-0.0859}
\expandafter\def\csname nhG@6@3\endcsname{-0.1178}
\expandafter\def\csname nhB@6@3\endcsname{-0.0024}
\expandafter\def\csname nhR@6@4\endcsname{0.0353}
\expandafter\def\csname nhG@6@4\endcsname{0.1228}
\expandafter\def\csname nhB@6@4\endcsname{-0.1255}
\expandafter\def\csname nhR@6@5\endcsname{-0.0938}
\expandafter\def\csname nhG@6@5\endcsname{0.0372}
\expandafter\def\csname nhB@6@5\endcsname{0.1103}
\expandafter\def\csname nhR@6@6\endcsname{-0.0181}
\expandafter\def\csname nhG@6@6\endcsname{0.0348}
\expandafter\def\csname nhB@6@6\endcsname{-0.1311}
\expandafter\def\csname nhR@6@7\endcsname{0.1372}
\expandafter\def\csname nhG@6@7\endcsname{-0.1189}
\expandafter\def\csname nhB@6@7\endcsname{0.0788}
\expandafter\def\csname nhR@7@0\endcsname{0.1239}
\expandafter\def\csname nhG@7@0\endcsname{0.0816}
\expandafter\def\csname nhB@7@0\endcsname{0.0713}
\expandafter\def\csname nhR@7@1\endcsname{0.0161}
\expandafter\def\csname nhG@7@1\endcsname{-0.0126}
\expandafter\def\csname nhB@7@1\endcsname{0.0633}
\expandafter\def\csname nhR@7@2\endcsname{-0.0998}
\expandafter\def\csname nhG@7@2\endcsname{-0.0862}
\expandafter\def\csname nhB@7@2\endcsname{0.0219}
\expandafter\def\csname nhR@7@3\endcsname{-0.0378}
\expandafter\def\csname nhG@7@3\endcsname{-0.1243}
\expandafter\def\csname nhB@7@3\endcsname{-0.0411}
\expandafter\def\csname nhR@7@4\endcsname{0.0108}
\expandafter\def\csname nhG@7@4\endcsname{-0.1465}
\expandafter\def\csname nhB@7@4\endcsname{0.0993}
\expandafter\def\csname nhR@7@5\endcsname{-0.0763}
\expandafter\def\csname nhG@7@5\endcsname{-0.0888}
\expandafter\def\csname nhB@7@5\endcsname{-0.1221}
\expandafter\def\csname nhR@7@6\endcsname{0.1064}
\expandafter\def\csname nhG@7@6\endcsname{0.1112}
\expandafter\def\csname nhB@7@6\endcsname{-0.0016}
\expandafter\def\csname nhR@7@7\endcsname{0.0646}
\expandafter\def\csname nhG@7@7\endcsname{-0.1211}
\expandafter\def\csname nhB@7@7\endcsname{0.0628}


\newcount\ugcnt   % active-node counter (logged via \message; should print 64/32/16/8)

% =====================================================================
\newcommand{\drawugraph}[5]{%
  % #1 = x-centre, #2 = y-centre, #3 = level (1..4), #4 = gid (unique string),
  % #5 = display mode: 0 clean | 1 noisy encoder | 2 clean estimate xhat | 3 noise estimate eps
  \begingroup
  \def\gid{#4}\def\noisemode{#5}%
  \global\ugcnt=0\relax
  % ---- per-level smoothing factor (smf) + selection threshold (thr) --
  \ifcase\numexpr#3-1\relax
      \def\smf{1.00}\def\Nlev{64}%   level 1 : finest, all 64 active
  \or \def\smf{0.62}\def\Nlev{32}%   level 2 : strict down-sample x2 -> 32
  \or \def\smf{0.42}\def\Nlev{16}%   level 3 : strict down-sample x2 -> 16
  \or \def\smf{0.28}\def\Nlev{8}%    level 4 : strict down-sample x2 -> 8
  \fi
  \pgfmathsetmacro{\gDlev}{\gDbase*\smf}%
  % ---- bump centres (local coords, relative to graph centre) ---------
  \pgfmathsetmacro{\bAx}{-0.42*\half}\pgfmathsetmacro{\bAy}{ 0.40*\half}%
  \pgfmathsetmacro{\bBx}{ 0.48*\half}\pgfmathsetmacro{\bBy}{ 0.06*\half}%
  \pgfmathsetmacro{\bCx}{-0.06*\half}\pgfmathsetmacro{\bCy}{-0.48*\half}%
  % ================= LOOP 1 : compute + store (no drawing) ============
  \foreach \i in {0,...,\gNm}{%
    \foreach \j in {0,...,\gNm}{%
      \pgfmathsetseed{\numexpr \gSeed + 7919*\i + 1231*\j \relax}%
      \pgfmathsetmacro{\lx}{(\i-\gC)*\gSp + \gJit*\gSp*rand}%
      \pgfmathsetmacro{\ly}{(\j-\gC)*\gSp + \gJit*\gSp*rand}%
      % normalised distances to the three bumps
      \pgfmathsetmacro{\rA}{sqrt((\lx-\bAx)^2+(\ly-\bAy)^2)/\half}%
      \pgfmathsetmacro{\rB}{sqrt((\lx-\bBx)^2+(\ly-\bBy)^2)/\half}%
      \pgfmathsetmacro{\rC}{sqrt((\lx-\bCx)^2+(\ly-\bCy)^2)/\half}%
      % SELECTION : nested top-N priority (baked perturbed energy+uniform rule)
      \edef\pri{\csname ugP@\i @\j\endcsname}%
      \pgfmathtruncatemacro{\acti}{\pri<\Nlev ? 1 : 0}%
      \ifnum\acti=1\relax\global\advance\ugcnt by1\relax\fi
      % SIGNAL colour : level-decay weighted mix of the bump colours (low-pass)
      \pgfmathsetmacro{\wA}{exp(-\gDlev*\rA)}%
      \pgfmathsetmacro{\wB}{exp(-\gDlev*\rB)}%
      \pgfmathsetmacro{\wC}{exp(-\gDlev*\rC)}%
      \pgfmathsetmacro{\sW}{max(\wA+\wB+\wC,1e-6)}%
      \pgfmathsetmacro{\Rc}{(\wA*\ugAR+\wB*\ugBR+\wC*\ugCR)/\sW}%
      \pgfmathsetmacro{\Gc}{(\wA*\ugAG+\wB*\ugBG+\wC*\ugCG)/\sW}%
      \pgfmathsetmacro{\Bc}{(\wA*\ugAB+\wB*\ugBB+\wC*\ugCB)/\sW}%
      % display mode: 0 clean field | 1 noisy encoder (clean+nz) | 2 clean estimate xhat (clean+nh)
      %               | 3 noise estimate eps=x_k-xhat (grey 0.5 + ne).  nz/ne are per-depth, nh full-res.
      \ifcase\noisemode\relax
        % mode 0 -- clean: leave \Rc,\Gc,\Bc as the smoothed clean field
      \or % mode 1 -- encoder: clean + LP residual noise
        \edef\dR{\csname nzR@#3@\i @\j\endcsname}\edef\dG{\csname nzG@#3@\i @\j\endcsname}\edef\dB{\csname nzB@#3@\i @\j\endcsname}%
        \pgfmathsetmacro{\Rc}{max(0,min(1,(\Rc)+(\dR)))}\pgfmathsetmacro{\Gc}{max(0,min(1,(\Gc)+(\dG)))}\pgfmathsetmacro{\Bc}{max(0,min(1,(\Bc)+(\dB)))}%
      \or % mode 2 -- clean estimate xhat_0: clean + small residual
        \edef\dR{\csname nhR@\i @\j\endcsname}\edef\dG{\csname nhG@\i @\j\endcsname}\edef\dB{\csname nhB@\i @\j\endcsname}%
        \pgfmathsetmacro{\Rc}{max(0,min(1,(\Rc)+(\dR)))}\pgfmathsetmacro{\Gc}{max(0,min(1,(\Gc)+(\dG)))}\pgfmathsetmacro{\Bc}{max(0,min(1,(\Bc)+(\dB)))}%
      \or % mode 3 -- noise estimate eps = x_k - xhat: grey base + LP eps (clean field cancels)
        \edef\dR{\csname neR@#3@\i @\j\endcsname}\edef\dG{\csname neG@#3@\i @\j\endcsname}\edef\dB{\csname neB@#3@\i @\j\endcsname}%
        \pgfmathsetmacro{\Rc}{max(0,min(1,0.5+(\dR)))}\pgfmathsetmacro{\Gc}{max(0,min(1,0.5+(\dG)))}\pgfmathsetmacro{\Bc}{max(0,min(1,0.5+(\dB)))}%
      \fi
      \pgfmathsetmacro{\gx}{#1+\lx}\pgfmathsetmacro{\gy}{#2+\ly}%
      \global\expandafter\edef\csname ugx@\gid @\i @\j\endcsname{\gx}%
      \global\expandafter\edef\csname ugy@\gid @\i @\j\endcsname{\gy}%
      \global\expandafter\edef\csname uga@\gid @\i @\j\endcsname{\acti}%
      \global\expandafter\edef\csname ugr@\gid @\i @\j\endcsname{\Rc}%
      \global\expandafter\edef\csname ugg@\gid @\i @\j\endcsname{\Gc}%
      \global\expandafter\edef\csname ugb@\gid @\i @\j\endcsname{\Bc}%
    }%
  }%
  \message{[ugraph \gid\space L#3] active=\the\ugcnt/\numexpr\gN*\gN\relax}%
  % ================= LOOP 2 : edges (active-active stronger) ==========
  \foreach \i in {0,...,\gNmm}{\foreach \j in {0,...,\gNm}{%   horizontal
    \pgfmathtruncatemacro{\ip}{\i+1}%
    \edef\xa{\csname ugx@\gid @\i @\j\endcsname}\edef\ya{\csname ugy@\gid @\i @\j\endcsname}%
    \edef\xb{\csname ugx@\gid @\ip @\j\endcsname}\edef\yb{\csname ugy@\gid @\ip @\j\endcsname}%
    \edef\aa{\csname uga@\gid @\i @\j\endcsname}\edef\ab{\csname uga@\gid @\ip @\j\endcsname}%
    \pgfmathtruncatemacro{\both}{\aa*\ab}%
    \ifnum\both=1\relax \draw[black!55,line width=0.40pt] (\xa,\ya)--(\xb,\yb);
    \else \draw[black!30,line width=0.30pt] (\xa,\ya)--(\xb,\yb);\fi
  }}%
  \foreach \i in {0,...,\gNm}{\foreach \j in {0,...,\gNmm}{%   vertical
    \pgfmathtruncatemacro{\jp}{\j+1}%
    \edef\xa{\csname ugx@\gid @\i @\j\endcsname}\edef\ya{\csname ugy@\gid @\i @\j\endcsname}%
    \edef\xb{\csname ugx@\gid @\i @\jp\endcsname}\edef\yb{\csname ugy@\gid @\i @\jp\endcsname}%
    \edef\aa{\csname uga@\gid @\i @\j\endcsname}\edef\ab{\csname uga@\gid @\i @\jp\endcsname}%
    \pgfmathtruncatemacro{\both}{\aa*\ab}%
    \ifnum\both=1\relax \draw[black!55,line width=0.40pt] (\xa,\ya)--(\xb,\yb);
    \else \draw[black!30,line width=0.30pt] (\xa,\ya)--(\xb,\yb);\fi
  }}%
  % ================= LOOP 3 : nodes ===================================
  \foreach \i in {0,...,\gNm}{\foreach \j in {0,...,\gNm}{%
    \edef\nx{\csname ugx@\gid @\i @\j\endcsname}\edef\ny{\csname ugy@\gid @\i @\j\endcsname}%
    \edef\na{\csname uga@\gid @\i @\j\endcsname}%
    \edef\nr{\csname ugr@\gid @\i @\j\endcsname}\edef\ng{\csname ugg@\gid @\i @\j\endcsname}\edef\nb{\csname ugb@\gid @\i @\j\endcsname}%
    \ifnum\na=1\relax
      \fill[draw=black!55,line width=0.2pt,fill={rgb,1:red,\nr;green,\ng;blue,\nb}] (\nx,\ny) circle (\gRa);
    \else
      \draw[black!28,line width=0.16pt,fill=white] (\nx,\ny) circle (\gRi);
    \fi
  }}%
  \endgroup
}
